% concatenated from main.tex
\documentclass[12pt]{amsart}
\usepackage[top=35mm, bottom=32mm, left=30mm, right=30mm]{geometry}
\usepackage{amsmath,amssymb,amsthm,mathtools}
\usepackage{enumitem}
\usepackage{microtype}
\usepackage[hidelinks]{hyperref}
\usepackage{mathrsfs}
\usepackage{bm}

\usepackage{aliascnt}

\newtheorem{theorem}{Theorem}[section]
\newaliascnt{proposition}{theorem}
\newtheorem{proposition}[proposition]{Proposition}
\aliascntresetthe{proposition}
\newaliascnt{lemma}{theorem}
\newtheorem{lemma}[lemma]{Lemma}
\aliascntresetthe{lemma}
\newaliascnt{corollary}{theorem}
\newtheorem{corollary}[corollary]{Corollary}
\aliascntresetthe{corollary}
\newaliascnt{claim}{theorem}
\newtheorem{claim}[claim]{Claim}
\aliascntresetthe{claim}
\theoremstyle{definition}
\newaliascnt{definition}{theorem}
\newtheorem{definition}[definition]{Definition}
\aliascntresetthe{definition}
\newaliascnt{remark}{theorem}
\newtheorem{remark}[remark]{Remark}
\aliascntresetthe{remark}

\usepackage[nameinlink,capitalise]{cleveref}
\crefname{theorem}{Theorem}{Theorems}
\crefname{proposition}{Proposition}{Propositions}
\crefname{lemma}{Lemma}{Lemmas}
\crefname{corollary}{Corollary}{Corollaries}
\crefname{claim}{Claim}{Claims}
\crefname{definition}{Definition}{Definitions}
\crefname{remark}{Remark}{Remarks}

\newcommand{\Haar}{m_G}
\newcommand{\vol}{\operatorname{vol}}
\newcommand{\Prob}{\operatorname{Prob}}
\newcommand{\Map}{\operatorname{Map}}
\newcommand{\Mapavg}{\operatorname{Map}_{\mathrm{avg}}}
\newcommand{\Sep}{\operatorname{Sep}}
\newcommand{\PPP}{\operatorname{PPP}}
\newcommand{\cN}{\mathcal N}
\newcommand{\cB}{\mathcal B}
\newcommand{\cO}{\mathcal O}
\newcommand{\1}{\mathbf 1}
\newcommand{\eps}{\varepsilon}
\newcommand{\dbar}{\overline{\vol}}
\newcommand{\dd}{\,d}
\newcommand{\norm}[1]{\left\lVert#1\right\rVert}
\newcommand{\abs}[1]{\left\lvert#1\right\rvert}

\title[Poisson actions and completely positive sofic entropy]{Poisson actions of noncompact locally compact sofic groups have completely positive entropy}

\author{Zhuowei Liu}
\date{\today}
\subjclass[2020]{37A35; 37A15}

\address[Z. Liu]{School of Mathematics (Zhuhai), Sun Yat-sen University,
	Zhuhai, Guangdong, 519000, P.R. China}
\email{liuzhw55@mail2.sysu.edu.cn}
\keywords{Locally compact sofic groups; Poisson actions; completely positive entropy}
\begin{document}

\begin{abstract}
In this paper, we use completed root views to study measure sofic entropy of
Poisson actions of locally compact sofic groups. We prove that, for every
noncompact locally compact second countable sofic group and every locally
compact sofic approximation, each positive-intensity Poisson action has
completely positive measure sofic entropy. 
We also compare the averaged and pointwise spatial formulas with Singh's entropy
and show that noncompactness is necessary for the theorem as stated.
\end{abstract}

\maketitle

\section{Introduction}

Completely positive entropy is stronger than positive entropy.
A probability-measure-preserving action has completely positive entropy when
each nontrivial factor has positive entropy. So this condition rules out
zero-entropy structure both in the original action and in all its factors.
In the classical case of a single transformation, completely
positive entropy is equivalent, by the Rokhlin--Sinai theorem, to the
$K$-property and is closely tied to strong mixing \cite{Glasner}. For countable amenable group actions, it also has strong mixing consequences
\cite{RudolphWeiss}. Bernoulli actions are the basic examples. The
Ornstein--Weiss theory shows that factors of Bernoulli actions of countably
infinite amenable groups are again Bernoulli, and so such actions have
completely positive entropy \cite{OrnsteinWeiss}.

Bowen's sofic entropy extends measure entropy beyond amenable groups by
counting finite models associated with a sofic approximation
\cite{BowenSofic,B12,B20,KerrLi1}. Kerr and Li developed an operator-algebraic approach to  Bowen’s
sofic measure entropy \cite{KerrLi2}. In this setting, the entropy depends on the chosen approximation. The
amenable theory of Bernoulli factors cannot be used to prove complete
positivity. But Kerr proved that every Bernoulli action of a countable sofic group has completely positive
entropy with respect to every sofic approximation \cite{KerrCPE}. The main
point of his argument is a positive lower bound for the local entropy of a
nontrivial partition which is uniform over all sufficiently good finite
models. This shows that complete positivity can be proved directly at the
level of arbitrary factors, without first showing that those factors are Bernoulli.
Austin and Burton later studied how this result is related to uniform mixing \cite{AustinBurton}.

Entropy theory for locally compact second countable groups requires a
different kind of finite model. Singh introduced an operator-algebraic theory
of locally compact sofic entropy and proved a spatial formula on compact
models \cite[Chapter~3]{SinghThesis}. Bowen and Burton later defined local
compact soficity using finite-volume local $G$-spaces
\cite{BowenBurton}, and Bowen developed the corresponding spatial entropy
theory, including measure-conjugacy invariance and the variational principle
\cite{BowenLCEntropy}. Instead of maps from a finite set labelled by group
elements, one now studies approximately equivariant maps from a finite-volume
space on which the group law is defined only locally. This local geometry is
exactly what makes the Poisson case both natural and technically different
from the discrete Bernoulli case.

Let $G$ be a locally compact sofic group and
$\Haar$ a Haar measure on $G$. Let $\cN(G)$ be the space of locally finite
integer-valued Radon measures on $G$, and, for $\alpha>0$, let
\[
 \mu_\alpha:=\PPP(\alpha\Haar)
\]
be the law of the Poisson point process of intensity $\alpha\Haar$. We use the
left action by right shifts
\[
 (g\omega)(A)=\omega(Ag),\qquad g\in G.
\]


Poisson actions are the continuous-group versions of Bernoulli shifts:
independence is indexed by disjoint regions rather than by individual group
coordinates. This analogy is exact for a countable discrete group equipped
with counting Haar measure. In fact,
\[
\omega\longmapsto \bigl(\omega(\{g\})\bigr)_{g\in G}
\]
identifies $(\cN(G),\mu_\alpha)$ with the Bernoulli shift having base law
$\operatorname{Pois}(\alpha)$. So Kerr's theorem settles the discrete case.
But for nondiscrete groups, the Poisson process has no
coordinatewise product presentation and must be modeled through local charts.

Singh proved that positive-intensity Poisson actions of nondiscrete locally
compact sofic groups have infinite measure sofic entropy
\cite[Chapter~5]{SinghThesis}. This computation concerns the full action.
Complete positivity is stronger: even an action of infinite
entropy may have a nontrivial zero-entropy factor. Because of Kerr's theorem
and because Poisson actions are the locally compact analogue of Bernoulli
shifts, Bowen asked whether Poisson actions have completely positive
entropy relative to every locally compact sofic approximation
\cite[Question~4]{BowenLCEntropy}.

In this paper, we answer Bowen's question for every noncompact locally
compact second countable sofic group. The main obstacle is that pulling a
Poisson process back through a local chart produces a truncated
configuration lying in a null set for the target Poisson law, where an
arbitrary measurable factor map is not controlled. We solve this problem by
adding an independent exterior Poisson process. The \emph{completed root view} then
has exactly the target Poisson law at every root
and can be passed through any measurable factor map. The following is the main result.

\begin{theorem}\label{thm:main}
Let $G$ be a noncompact locally compact second countable sofic group, and
let $\Sigma=(M_i)_{i\geq1}$ be a locally compact sofic approximation of
$G$. For $\alpha>0$, let $\mu_\alpha:=\PPP(\alpha\Haar)$
be the law of the Poisson point process on $G$ with intensity measure
$\alpha\Haar$. Then the Poisson action
	\[
	G\curvearrowright(\cN(G),\mu_\alpha),
	\qquad
	(g\omega)(A)=\omega(Ag),
	\]
	has completely positive measure sofic entropy with respect to $\Sigma$.
	In other words, every nontrivial measurable factor
	\[
	\pi:(\cN(G),\mu_\alpha)\longrightarrow(Z,\zeta)
	\]
	satisfies
	\[
	h^{\mathrm{meas}}_\Sigma
	\bigl(G\curvearrowright(Z,\zeta)\bigr)>0.
	\]
\end{theorem}

The exact law of the completed root views gives well-defined random models
on every measurable factor. Local covariance and Poisson concentration,
followed by a packing argument, produce exponentially many separated factor
microstates. We also prove the equivalence of the averaged and pointwise
spatial entropy formulas and compare them with Singh's entropy.
Noncompactness is essential: it gives
\[
\vol_{M_i}(M_i)\longrightarrow\infty,
\]
while for a nontrivial compact group the total number of Poisson points
defines a nontrivial invariant factor.

\medskip
\noindent\textbf{Organization of the paper.}
Section~2 recalls the spatial entropy formalism and proves the equivalence of
averaged and pointwise equivariance. Section~3 collects the geometric
properties of locally compact sofic models. Section~4 constructs the completed
root views and proves their exact marginal and local covariance properties.
Section~5 proves exponential concentration for local Poisson observables.
Sections~6--7 pass the construction to arbitrary measurable factors and prove
the separation. Section~8 proves  packing estimates and  the main theorem.
Section~9 gives the compact-group obstruction. The final section proves the
invariance statements used above and compares the spatial formula with Singh's
entropy.


\section{Spatial measure sofic entropy on compact models}\label{sec:compact-entropy}




\subsection{Standard probability-measure-preserving actions}
\begin{definition}[Standard pmp action]
	Let $G$ be a locally compact second countable group.
	A standard probability-measure-preserving (pmp) action of $G$ is a
	Borel action
	\[
	G\times X\to X,\qquad (g,x)\mapsto gx,
	\]
	on a standard probability space $(X,\mu)$ such that
$\mu(gA)=\mu(A)$
	for every $g\in G$ and every Borel set $A\subset X$.	
	In other words, the probability measure $\mu$ is invariant under the
	action of $G$.
\end{definition}
We use the abbreviation lcsc to mean locally compact second
countable.
\begin{definition}[Measurable factor]
	Let
	$G\curvearrowright(X,\mu)$,
	$G\curvearrowright(Y,\nu)$
	be standard pmp actions. A measurable \emph{factor map} is a measurable map
	$\pi:X\to Y$
	such that
$\pi_*\mu=\nu$ (where $\pi_*\mu$ is the measure on
$Y$ defined by $\pi_*\mu(A)=\mu (\pi^{-1}A)$)
	and
$\pi(gx)=g\pi(x)$
	for every $g\in G$ and for $\mu$-almost every $x\in X$; a \emph{measure conjugacy} is a factor map for which there exist invariant
	conull Borel sets $X_0\subseteq X$ and $Y_0\subseteq Y$ such that
	$\pi|_{X_0}:X_0\to Y_0$ is a Borel bijection with Borel inverse and
	intertwines the two actions on these conull sets.

\end{definition}

\begin{definition}
A \emph{compact model} of a standard pmp action $G\curvearrowright(Z,\zeta)$
consists of a compact metrizable space $Y$, a jointly continuous action
$G\curvearrowright Y$, an invariant Borel probability measure $\nu$, and a
measure conjugacy between $(Z,\zeta)$ and $(Y,\nu)$.
\end{definition}

\begin{proposition}[Existence of compact models]\label{prop:compact-model}\cite{Var63}
Every standard pmp action of an lcsc group admits a compact metrizable
continuous model.
\end{proposition}

We work throughout with compact metrizable target systems. The Poisson
configuration space itself is noncompact and appears only as the source of the
factor map and as an auxiliary space in the random construction.


\subsection{Local \texorpdfstring{$G$}{G}-spaces and sofic approximations}

We recall the local-$G$-space formalism of Bowen and Burton
\cite{BowenBurton}. Let G be an lcsc group with identity $e$.

\begin{definition}\label{defn}
Let $M$ be a locally compact second countable Hausdorff space. A map
$\vartheta:\operatorname{dom}(\vartheta)\to M$, where
$\operatorname{dom}(\vartheta)\subseteq G\times M$ is open, is a
\emph{partial left action} of $G$ if, for every $p\in M$,
\begin{enumerate}
\item $(e,p)\in\operatorname{dom}(\vartheta)$ and
$\vartheta(e,p)=p$;
\item if $(g,p)\in\operatorname{dom}(\vartheta)$, then
$(g^{-1},\vartheta(g,p))\in\operatorname{dom}(\vartheta)$ and
$\vartheta(g^{-1},\vartheta(g,p))=p$;
\item if $(h,p)$, $(g,\vartheta(h,p))$, and $(gh,p)$ belong to
$\operatorname{dom}(\vartheta)$, then
$\vartheta(gh,p)=\vartheta(g,\vartheta(h,p))$.
\end{enumerate}
The partial action is \emph{homogeneous} if, in addition,
\begin{enumerate}[resume]
\item for every $p\in M$, there is an open neighborhood $V_p$ of $e$ such
that $V_p\times\{p\}\subseteq\operatorname{dom}(\vartheta)$ and
$\vartheta(\cdot,p)$ restricts to a homeomorphism from $V_p$ onto an open
neighborhood of $p$ in $M$.
\end{enumerate}
A pair $(M,\vartheta)$ is a \emph{local left $G$-space} if $\vartheta$ is a
homogeneous partial left action of $G$ on $M$.
\end{definition}

We usually suppress $\vartheta$ and write $g.p$ for $\vartheta(g,p)$. If
$L\subseteq M$ and $\{g\}\times L\subseteq\operatorname{dom}(\vartheta)$, we
write $g.L:=\{g.p:p\in L\}$. In the same way, if $V\subseteq G$ and
$V\times\{p\}\subseteq\operatorname{dom}(\vartheta)$, we write
$V.p:=\{g.p:g\in V\}$.

\begin{definition}\label{def:chart}
Let $(M,\vartheta)$ be a local $G$-space and let $p\in M$. A
\emph{chart centered at $p$} is a homeomorphism
$f_p:\operatorname{dom}(f_p)\to\operatorname{rng}(f_p)$ such that
\begin{enumerate}
\item $\operatorname{dom}(f_p)$ is an open neighborhood of $p$ in $M$;
\item $\operatorname{rng}(f_p)$ is an open neighborhood of $e$ in $G$;
\item $f_p(g.p)=g$ for every $g\in\operatorname{rng}(f_p)$.
\end{enumerate}
Condition~(4) of \cref{defn} gives such a chart at every point.
\end{definition}

The next two results from \cite{BowenBurton} provide the canonical measure on
a local $G$-space and its local invariance.

\begin{proposition}
Let $(M,\vartheta)$ be a local $G$-space, and fix a right Haar measure $\Haar$
on $G$. There is a unique Radon measure $\vol_M$ on $M$ such that, whenever
$f_p$ is a chart centered at $p$ and $L\subseteq\operatorname{dom}(f_p)$ is
Borel,
\[
 \vol_M(L)=\Haar\bigl(\{g\in\operatorname{rng}(f_p):g.p\in L\}\bigr)
 =\Haar(f_p(L)).
\]
When the underlying local $G$-space is clear, we write $\vol$ for $\vol_M$.
\end{proposition}

\begin{lemma}\label{lmp}
Let $(M,\vartheta)$ be a local $G$-space, and suppose that
$\{g\}\times L\subseteq\operatorname{dom}(\vartheta)$ for some measurable
$L\subseteq M$ and $g\in G$. If $G$ is unimodular, then
\[
 \vol(g.L)=\vol(L).
\]
\end{lemma}

\begin{definition}\label{def:local-sofic-model}
Let $\mathcal M=(M,\vartheta)$ be a local $G$-space, let $U\subseteq G$ be
open and precompact, and let $\eps>0$. Define $M[U]=M[\vartheta,U]$ to be the
set of all $p\in M$ such that
\begin{enumerate}
\item if $g,h,gh\in U$, then $g.(h.p)=(gh).p$, with both sides defined;
\item the map $g\mapsto g.p$ is a homeomorphism from $U$ onto an open
neighborhood of $p$.
\end{enumerate}
If $0<\vol(M)<\infty$ and
\[
 \vol(M[U])\geq(1-\eps)\vol(M),
\]
then $\mathcal M$ is a \emph{$(U,\eps)$-sofic approximation} to $G$.
\end{definition}

\begin{definition}\label{def:sofic-approximation}
A \emph{sofic approximation} to $G$ is a sequence
$\Sigma=(M_i)_{i\geq1}$ such that each $M_i$ is a
$(U_i,\eps_i)$-sofic approximation, where
\begin{enumerate}
\item the $U_i$ are precompact open sets satisfying
$U_1\subseteq U_2\subseteq\cdots$ and $\bigcup_iU_i=G$;
\item $\eps_i\to0$ as $i\to\infty$.
\end{enumerate}
The group $G$ is \emph{sofic} if it admits such an approximation.
\end{definition}
% For a sofic approximation $\Sigma=(M_i)$, we write
%$V_i:=\operatorname{vol}_{M_i}(M_i)$
%for the volume of the $i$-th local model and $\dbar_i:=V_i^{-1}\vol_{M_i}$.
\subsection{The spatial microstate formula}

 We use Bowen's spatial definition of measure sofic entropy for locally compact
 groups, specialized to compact metrizable continuous pmp targets. It is the
 local-$G$-space version of Singh's pseudometric formula
 \cite[Section~3.4]{SinghThesis}. Singh uses a normalized $L^2$ metric and
 finitely many continuous test functions, while Bowen uses a partial-map
 quasimetric and weak-star neighborhoods of the invariant measure. The two
 forms of the empirical-law condition are equivalent, and the comparison with
 Singh's normalized $L^2$ formulation is proved in
 \cref{app:singh-comparison}.
	We include metric independence and topological-conjugacy invariance in
\cref{app:singh-comparison}. The stronger measure-conjugacy invariance and
compact-model independence used in the main theorem follow from the explicit
comparison with Singh's measure entropy.

Let $G\curvearrowright(Y,\rho)$ be a jointly continuous action on a compact metric space, with $0\leq\rho\leq1$, and let $\nu$ be $G$-invariant. Fix a proper left-invariant compatible metric $d_G$ on $G$ and write
\[
 B(r):=\{g\in G:d_G(g,e)<r\}.
\]
If $M$ is a finite-volume local $G$-space, write
\[
 \overline{\vol}_M:=\frac{\vol_M}{\vol_M(M)},
 \qquad
 \rho^M(\psi,\psi'):=\int_M\rho(\psi(p),\psi'(p))\,d\overline{\vol}_M(p).
\]
For $g\in G$ let
\[
 M[g]:=\{p\in M:(g,p)\text{ lies in the domain of the partial action}\}.
\]
For a measurable map $\psi:M\to Y$ define
\begin{equation}\label{eq:defect}
 \Delta_{\rho,M}(g,\psi)
 :=\int_{M[g]}\rho\bigl(\psi(g.p),g\psi(p)\bigr)\,d\overline{\vol}_M(p)
   +\overline{\vol}_M(M\setminus M[g]).
\end{equation}
Undefined points get the maximal penalty $1$. In other words, this is Bowen's quasi-metric distance between the partially defined map $\psi\circ g$ and the everywhere defined map $g\circ\psi$.

For a precompact open identity neighborhood $L\subset G$, $\delta>0$, and a weak-star neighborhood $\cO$ of $\nu$, define the standard model space
\[
 \Map(M,Y,\rho:L,\delta,\cO)
 :=\left\{\psi:M\to Y:
 \Delta_{\rho,M}(g,\psi)<\delta\ \forall g\in L,
 \quad \psi_*\overline{\vol}_M\in\cO\right\}.
\]
We also define the averaged model space
\[
 \Mapavg(M,Y,\rho:L,\delta,\cO)
 :=\left\{\psi:M\to Y:
 \begin{aligned}
 &\frac1{\Haar(L)}\int_L\Delta_{\rho,M}(g,\psi)\,d\Haar(g)<\delta,\\
 &\psi_*\overline{\vol}_M\in\cO
 \end{aligned}
 \right\}.
\]
When the empirical condition is omitted, we suppress $\cO$ from the notation.



\begin{definition}
Let $(Y,\rho)$ be a pseudometric space and let $\eps>0$. A subset
$E\subseteq Y$ is \emph{$(\rho,\eps)$-separated} if
$\rho(x,y)>\eps$ for all distinct $x,y\in E$. We write
$\Sep_\eps(E,\rho^M)$ for the
maximum cardinality of a
$(\rho^M,\eps)$-separated subset of $E$, with the conventions
$\Sep_\eps(\varnothing,\rho^M)=0$, $\log0=-\infty$ and
$\log(+\infty)=+\infty$.
\end{definition}

\begin{definition}\label{def:spatial-entropy}
For $\Sigma=(M_i)$,  $V_i:=\vol_{M_i}(M_i)$. The spatial measure sofic
entropy of the compact model $(Y,\nu)$ relative to $\Sigma$ is defined by
\begin{align*}
 h^{\mathrm{meas}}_{\Sigma,\eps}(G,Y,\nu,\rho)
 &:={\inf}_{\cO\ni\nu}\,{\inf}_{L}\,{\inf}_{\delta>0}
 \limsup_{i\to\infty}\frac1{V_i}
 \log\Sep_\eps\bigl(\Map(M_i,Y,\rho:L,\delta,\cO),\rho^{M_i}\bigr),\\
 h^{\mathrm{meas}}_{\Sigma}(G,Y,\nu,\rho)
 &:=\lim_{\eps\downarrow0}
 h^{\mathrm{meas}}_{\Sigma,\eps}(G,Y,\nu,\rho),
\end{align*}
where $L$ ranges over precompact open identity neighborhoods of $G$.
Replacing $\Map$ by $\Mapavg$ defines
$h^{\mathrm{meas,avg}}_{\Sigma,\eps}$ and
$h^{\mathrm{meas,avg}}_\Sigma$.
\end{definition}
\begin{remark}
	The appendix compares the local-$G$-space formula above with Singh's
compact-model pseudometric entropy for the completed approximation
$\widetilde\Sigma$.  When $G$ is nondiscrete, Singh's
Theorem~3.4.2, Proposition~3.4.3, and Definition~11 identify that
pseudometric entropy with his standard measure entropy.  Since the latter
is a measure-conjugacy invariant, the value is independent of the chosen
compact model.
\end{remark}

The next comparison connects the random construction with Bowen's pointwise spatial formula. It is the empirical-law version of \cite[Proposition~4.10]{BowenLCEntropy}; we include the proof because preservation of the weak-star condition is essential below.

\begin{lemma}\label{lem:large-products}\cite[Lemma~4.11]{BowenLCEntropy}
Let $0<r_0<r_1<r_2$ and $0<\tau<1$. Assume $r_0+r_1<r_2$ and
\begin{equation}\label{eq:haar-balls-condition}
 \tau\Haar(B(r_2))<(1-\tau)\Haar(B(r_0)).
\end{equation}
If measurable sets $W_j\subset B(r_j)$ satisfy
\[
 \Haar(W_j)>(1-\tau)\Haar(B(r_j)),\qquad j=0,2,
\]
then
\[
 W_0W_2\supset B(r_1).
\]
\end{lemma}



\begin{proposition}[Averaged and pointwise spatial formulas coincide]\label{prop:avg-standard}
For every compact metrizable continuous $G$-system $(Y,\nu)$ and every compatible metric $\rho\leq1$,
\[
 h^{\mathrm{meas,avg}}_{\Sigma,\eps}(G,Y,\nu,\rho)
 =h^{\mathrm{meas}}_{\Sigma,\eps}(G,Y,\nu,\rho)
 \qquad(\eps>0).
\]
So we have
\[
 h^{\mathrm{meas,avg}}_\Sigma(G,Y,\nu,\rho)
 =h^{\mathrm{meas}}_\Sigma(G,Y,\nu,\rho).
\]
\end{proposition}

\begin{proof}
Since every pointwise microstate is averaged,
\[
 \Map(M,Y,\rho:L,\delta,\cO)
 \subseteq \Mapavg(M,Y,\rho:L,\delta,\cO),
\]
and so
\[
 h^{\mathrm{meas}}_{\Sigma,\eps}(G,Y,\nu,\rho)
 \leq h^{\mathrm{meas,avg}}_{\Sigma,\eps}(G,Y,\nu,\rho).
\]

For the reverse inequality, we show that one model space is eventually contained in the other. Fix a precompact open identity neighborhood $L\subset G$, a number $0<\delta<1$, and a weak-star neighborhood $\cO$ of $\nu$. We will choose a precompact open identity neighborhood $\widetilde L$, a tolerance $\widetilde\delta>0$, and an index $I$ such that
\begin{equation}\label{eq:avg-to-standard-inclusion}
 \Mapavg(M_i,Y,\rho:\widetilde L,\widetilde\delta,\cO)
 \subseteq
 \Map(M_i,Y,\rho:L,\delta,\cO)
 \qquad(i\geq I).
\end{equation}
We do not change the microstate, so the empirical-measure condition is preserved exactly.

By joint continuity of the action and compactness of $Y$, the map
\[
 G\times Y\times Y\longrightarrow\mathbb R,
 \qquad (h,x,y)\longmapsto\rho(hx,hy),
\]
is uniformly continuous on a sufficiently small compact neighborhood of $\{e\}\times Y\times Y$. So there are $r_0>0$ and $\kappa>0$ such that
\begin{equation}\label{eq:local-equicontinuity}
 \rho(x,y)<\kappa,\quad h\in B(r_0)
 \quad\Longrightarrow\quad
 \rho(hx,hy)<\delta/10.
\end{equation}
Choose $r_1>r_0$ with $L\subset B(r_1)$, and then choose $r_2>r_0+r_1$. Since $B(r_0)$ and $B(r_2)$ have finite positive Haar measure, we may choose $0<\tau<1$ so small that
\begin{equation}\label{eq:tau-choices}
 \tau\Haar(B(r_2))<(1-\tau)\Haar(B(r_0)),
 \qquad
 \tau<\min\{\kappa,\delta/10\}.
\end{equation}
Put
\[
 \widetilde L:=B(r_2),\qquad \widetilde\delta:=\tau^{10}.
\]
The balls are symmetric because
$d_G(g^{-1},e)=d_G(e,g)=d_G(g,e)$.

Because $U_i\uparrow G$, for all sufficiently large $i$ one has
$\widetilde L^2\subset U_i$. Monotonicity of the good sets then gives
$M_i[U_i]\subset M_i[\widetilde L^2]$, so the defining sofic
condition implies
$\dbar_{M_i}(M_i[\widetilde L^2])\to1$. Choose $i$ sufficiently large
that
\begin{equation}\label{eq:comparison-core}
 \dbar_{M_i}\bigl(M_i[\widetilde L^2]\bigr)
 >1-\min\{\tau^{10},\delta/10\}.
\end{equation}
Let
\[
 \psi\in\Mapavg(M_i,Y,\rho:\widetilde L,\tau^{10},\cO).
\]
For $u\in\widetilde L$, define
\[
 E(u):=\left\{p\in M_i[u]:
 \rho\bigl(u\psi(p),\psi(u.p)\bigr)<\tau^4\right\}
\]
and
\[
 W:=\left\{u\in\widetilde L:
 \dbar_{M_i}(E(u))>1-\tau^4\right\}.
\]
If $u\notin W$, then the complement of $E(u)$, including points at which $u.p$ is undefined, has normalized measure at least $\tau^4$. On the part of $M_i[u]\setminus E(u)$ the metric defect is at least $\tau^4$, while on $M_i\setminus M_i[u]$ the undefined-domain penalty is $1\geq\tau^4$. This gives
\[
 \Delta_{\rho,M_i}(u,\psi)
 \geq \tau^4\bigl(1-\dbar_{M_i}(E(u))\bigr)
 \geq\tau^8.
\]
The averaged microstate inequality now gives
\begin{align*}
 \tau^{10}\Haar(\widetilde L)
 &>\int_{\widetilde L}\Delta_{\rho,M_i}(u,\psi)\,d\Haar(u)\\
 &\geq\tau^8\Haar(\widetilde L\setminus W).
\end{align*}
So we have
\begin{equation}\label{eq:good-group-elements}
 \Haar(\widetilde L\setminus W)<\tau^2\Haar(B(r_2)).
\end{equation}
In particular,
\[
 \Haar(W)>(1-\tau^2)\Haar(B(r_2))
 >(1-\tau)\Haar(B(r_2)).
\]
Also, using the first inequality in \eqref{eq:tau-choices},
\begin{align*}
 \Haar(B(r_0)\setminus W)
 &\leq \Haar(\widetilde L\setminus W)\\
 &<\tau^2\Haar(B(r_2))\\
 &<\tau(1-\tau)\Haar(B(r_0))
 <\tau\Haar(B(r_0)).
\end{align*}
So we have
\[
 \Haar(W\cap B(r_0))>(1-\tau)\Haar(B(r_0)).
\]
By \cref{lem:large-products}, every $g\in L\subset B(r_1)$ admits a factorization
\begin{equation}\label{eq:g-hk-factorization}
 g=hk,
 \qquad h\in W\cap B(r_0),\quad k\in W.
\end{equation}
Fix such a $g$ and such a factorization.

We now estimate the defect of $\psi$ at $g$ directly. For a point $p\in M_i$, define three nonnegative functions $A_1,A_2,A_3$, each bounded by $1$, as follows. If the expressions involved are defined, let
\begin{align*}
 A_1(p)&:=\rho\bigl(h(k\psi(p)),h\psi(k.p)\bigr),\\
 A_2(p)&:=\rho\bigl(h\psi(k.p),\psi(h.(k.p))\bigr),\\
 A_3(p)&:=\rho\bigl(\psi(h.(k.p)),\psi((hk).p)\bigr);
\end{align*}
and set the corresponding $A_j(p)$ equal to $1$ whenever one of its required partial-action expressions is undefined. If $(hk).p$ is defined and all intermediate expressions are defined, the triangle inequality gives
\[
 \rho\bigl((hk)\psi(p),\psi((hk).p)\bigr)
 \leq A_1(p)+A_2(p)+A_3(p).
\]
If $(hk).p$ is undefined, then $A_3(p)=1$, so the undefined-domain penalty in \eqref{eq:defect} is also bounded by $A_1(p)+A_2(p)+A_3(p)$. If $(hk).p$ is defined but one of the intermediate expressions is undefined, then by definition at least one of $A_1(p),A_2(p),A_3(p)$ equals $1$, while the metric defect is at most $1$. In every case, the same pointwise bound holds, so
\begin{equation}\label{eq:defect-three-direct}
 \Delta_{\rho,M_i}(g,\psi)
 \leq \int_{M_i}\bigl(A_1+A_2+A_3\bigr)\,d\dbar_{M_i}.
\end{equation}

Because $k\in W$, we have $\dbar_{M_i}(E(k))>1-\tau^4$. For $p\in E(k)$,
\[
 \rho\bigl(k\psi(p),\psi(k.p)\bigr)<\tau^4<\kappa.
\]
Since $h\in B(r_0)$, \eqref{eq:local-equicontinuity} gives $A_1(p)<\delta/10$. On the complement of $E(k)$ we only use $A_1\leq1$. This gives
\begin{equation}\label{eq:A1-bound}
 \int_{M_i}A_1\,d\dbar_{M_i}
 <\frac\delta{10}+\tau^4.
\end{equation}

To estimate $A_2$, we must use the fact that a partial translation preserves the canonical measure. Set
\[
 C_i:=M_i[\widetilde L^2],
 \qquad F:=E(h)\cap C_i.
\]
Since $h\in W$, \eqref{eq:comparison-core} gives
\begin{equation}\label{eq:F-lower-measure}
 \dbar_{M_i}(F)
 \geq1-\tau^4-\dbar_{M_i}(M_i\setminus C_i)
 >1-\tau^4-\delta/10.
\end{equation}
For every $q\in F$, the point $k^{-1}.q$ is defined because $q\in M_i[\widetilde L^2]$ and $k^{-1}\in\widetilde L$. Put
\[
 p:=k^{-1}.q.
\]
The local inverse and multiplication axioms on $M_i[\widetilde L^2]$ imply
\[
 k.p=k.(k^{-1}.q)=q.
\]
So we have
\[
 k^{-1}.F\subseteq\{p\in M_i:k.p\in E(h)\}.
\]
The partial translation $T(q):=k^{-1}.q$ is injective on $F$, with inverse $p\mapsto k.p$ on $T(F)$. Since $F\subset M_i[\widetilde L^2]$ and $k,k^{-1}\in\widetilde L$, \cref{lem:local-measure-preserving} applies and gives
\[
 \vol_{M_i}(k^{-1}.F)=\vol_{M_i}(F).
\]
Combining this identity with \eqref{eq:F-lower-measure}, we get
\begin{equation}\label{eq:pullback-Eh}
 \dbar_{M_i}\{p:k.p\in E(h)\}
 >1-\tau^4-\delta/10.
\end{equation}
Whenever $k.p\in E(h)$, all expressions in $A_2(p)$ are defined and
\[
 A_2(p)=\rho\bigl(h\psi(k.p),\psi(h.(k.p))\bigr)<\tau^4.
\]
Using $A_2\leq1$ elsewhere and \eqref{eq:pullback-Eh}, we get
\begin{equation}\label{eq:A2-bound}
 \int_{M_i}A_2\,d\dbar_{M_i}
 \leq \tau^4+\tau^4+\delta/10.
\end{equation}

Finally, if $p\in C_i=M_i[\widetilde L^2]$, then $h,k,hk\in\widetilde L^2$ and the local multiplication axiom gives
\[
 h.(k.p)=(hk).p.
\]
So $A_3(p)=0$ on $C_i$, and $A_3\leq1$ everywhere. By \eqref{eq:comparison-core},
\begin{equation}\label{eq:A3-bound}
 \int_{M_i}A_3\,d\dbar_{M_i}
 \leq\dbar_{M_i}(M_i\setminus C_i)<\delta/10.
\end{equation}

Substituting \eqref{eq:A1-bound}, \eqref{eq:A2-bound}, and \eqref{eq:A3-bound} into \eqref{eq:defect-three-direct} gives
\[
 \Delta_{\rho,M_i}(g,\psi)
 <\frac{3\delta}{10}+3\tau^4.
\]
Since $\tau<\delta/10<1$, we have $\tau^4<\delta/10$, and so
\[
 \Delta_{\rho,M_i}(g,\psi)<\frac{6\delta}{10}<\delta.
\]
The element $g\in L$ was arbitrary, so $\psi$ belongs to
$\Map(M_i,Y,\rho:L,\delta,\cO)$. This proves \eqref{eq:avg-to-standard-inclusion}.

For every $\eps>0$, the inclusion gives
\begin{align*}
 &\limsup_{i\to\infty}\frac1{V_i}\log\Sep_\eps
 \bigl(\Mapavg(M_i,Y,\rho:\widetilde L,\widetilde\delta,\cO),\rho^{M_i}\bigr)\\
 &\qquad\leq
 \limsup_{i\to\infty}\frac1{V_i}\log\Sep_\eps
 \bigl(\Map(M_i,Y,\rho:L,\delta,\cO),\rho^{M_i}\bigr).
\end{align*}
Taking the infimum over the averaged parameters on the left and then over $L$, $\delta$, and $\cO$ on the right proves
\[
 h^{\mathrm{meas,avg}}_{\Sigma,\eps}
 \leq h^{\mathrm{meas}}_{\Sigma,\eps}.
\]
Together with the first inequality, this proves equality at every scale $\eps$, and taking the limit as $\eps\downarrow0$ proves the equality of the full entropies.
\end{proof}



%\begin{proposition}[Identification with the Singh--Bowen measure entropy]
%\label{thm:identification-bowen}
%For every compact metrizable continuous pmp $G$-system $(Y,\nu)$ and every
%compatible metric $\rho\leq1$, the quantity
%\[
% h^{\mathrm{meas}}_\Sigma(G,Y,\nu,\rho)
%\]
%defined above is the standard locally compact measure sofic entropy. It is
%the $L^1$ version of Singh's compact-model pseudometric formula and the
%compact-target specialization of Bowen's local-$G$-space formula. In
%particular, it is invariant under measure conjugacy and is independent of the
%chosen compact model of the underlying standard pmp action.
%\end{proposition}


	\begin{definition}\label{def:cpe}
	A standard pmp action $G\curvearrowright(X,\mu)$ has
	\emph{completely positive measure sofic entropy relative to $\Sigma$} if
	every nontrivial measurable factor $G\curvearrowright(Z,\zeta)$ has
	strictly positive standard measure sofic entropy relative to $\Sigma$.
	When $G$ is nondiscrete, \cref{cor:singh-measure-identification} shows that
	this entropy may be computed from any compact metrizable continuous model
	$(Y,\nu)$ of the factor as
	\[
	h^{\mathrm{meas}}_\Sigma(G,Y,\nu,\rho),
	\]
	where $\rho\leq1$ is any compatible metric. The resulting value is
	independent of both the compact model and the metric.
\end{definition}
\begin{remark}\label{rem:cpe-model-independence}
	This compact-model independence comes from the Singh--Bowen comparison and
	Singh's measure-conjugacy invariance; it does not follow just from
	topological-conjugacy invariance of the compact-model formula.
\end{remark}

%\begin{remark}\label{rem:notation-entropy}
%In the proof we represent a measurable factor by an arbitrary compact
%metrizable continuous model, whose existence is guaranteed by
%\cref{prop:compact-model}. By \cref{thm:identification-bowen}, the entropy of
%that compact model is the standard measure sofic entropy of the factor and is
%independent of the model.
%\end{remark}

\section{Locally compact sofic models}

Throughout the remainder of the paper, $G$ denotes the locally compact second countable sofic group in \cref{thm:main}. A local $G$-space, its canonical measure, and the good sets $M[U]$ are as in Section~2. The good sets $M[U]$ are Borel in the standard formulation. Fix a right Haar measure $\Haar$ on $G$.
Since locally compact sofic groups are unimodular
\cite{BowenBurton}, $\Haar$ is also left invariant and invariant under
inversion.

Let $\Sigma=(M_i)$ be a sofic approximation. So there are precompact open identity neighborhoods $U_i\uparrow G$ and numbers $\eps_i\downarrow0$ such that
\[
 \vol_{M_i}(M_i[U_i])\geq (1-\eps_i)\vol_{M_i}(M_i).
\]
Write
\[
V_i:=\vol_{M_i}(M_i),\qquad \dbar_i:=V_i^{-1}\vol_{M_i}.
\]



\begin{lemma}[Local translations preserve canonical measure]\label{lem:local-measure-preserving}
Let $M$ be a local $G$-space with canonical measure.
For every $g\in G$, the partial translation
\[
T_g:M[g]\to M[g^{-1}],\qquad T_g(p)=g.p,
\]
is a Borel bijection with inverse $T_{g^{-1}}$, and for every Borel
$E\subset M[g]$ one has
\[
\operatorname{vol}_M(T_gE)=\operatorname{vol}_M(E).
\]
\end{lemma}


\begin{proof}
	The inverse and multiplication axioms for a local action give
	$T_{g^{-1}}T_g=\mathrm{id}$ and $T_gT_{g^{-1}}=\mathrm{id}$ wherever the
	expressions are defined. So it is enough to prove measure
	preservation. Around every point of $M[g]$ there is a local orbit chart on
	which $T_g$ is represented in group coordinates by left multiplication by
	$g$. In these coordinates the canonical measure is $\Haar$, which is left
	invariant because $G$ is unimodular. Cover a Borel set $E\subset M[g]$ by
	countably many such chart domains and refine the cover to a disjoint Borel
	partition. The images of the pieces are disjoint because $T_g$ is injective.
	Chartwise invariance and countable additivity give
	$\vol_M(T_gE)=\vol_M(E)$. Borelness of $T_gE$ follows from Lusin--Souslin's theorem.
\end{proof}








\subsection{Divergence of model volume}

\begin{lemma}\label{lem:volume-divergence}
If $G$ is noncompact, then $V_i\to\infty$.
\end{lemma}

\begin{proof}
For all sufficiently large $i$, $\eps_i<1$, and so $M_i[U_i]$ is nonempty. Choose $p_i\in M_i[U_i]$. By the definition of $M_i[U_i]$, the orbit map
\[
 U_i\longrightarrow U_i.p_i,\qquad g\longmapsto g.p_i,
\]
is a homeomorphism onto an open neighborhood of $p_i$. By the defining property of the canonical measure,
\[
 \vol_{M_i}(U_i.p_i)=\Haar(U_i).
\]
So $V_i\geq \Haar(U_i)$. Since $U_i\uparrow G$, continuity from below gives
\[
 \lim_i\Haar(U_i)=\Haar(G).
\]
A locally compact group has finite Haar measure if and only if it is compact. So $\Haar(G)=\infty$, and the result follows.
\end{proof}

\subsection{Slowly growing chart windows}

Choose symmetric precompact open neighborhoods
\[
 B_1\subset B_2\subset\cdots,\qquad \overline{B_n}\subset B_{n+1},\qquad \bigcup_nB_n=G,
\]
with $e\in B_1$. By passing to a slowly increasing integer sequence $r(i)\to\infty$, we may arrange that
\[
 W_i:=B_{r(i)}\quad\text{satisfies}\quad W_i^4\subset U_i.
\]
Define
\[
 A_i:=M_i[W_i^2],\qquad D_i:=M_i[W_i^4].
\]
Since $W_i^4\subset U_i$, the monotonicity of the good sets gives
\[
 M_i[U_i]\subset D_i\subset A_i.
\]
So we have
\begin{equation}\label{eq:deep-core-density}
 d_i:=\dbar_i(D_i)\geq1-\eps_i\longrightarrow1.
\end{equation}

The following change-of-root lemma explains the action convention used throughout the proof.

\begin{lemma}[Change of root]\label{lem:change-root}
Let $K,L\subset G$ be compact. For all sufficiently large $i$, the following holds. If $p\in D_i$, $g\in L$, and $q=g.p$, then $q\in A_i$ and
\[
 h.q=(hg).p\qquad\text{for every }h\in K.
\]
Also, the orbit maps $h\mapsto h.p$ and $h\mapsto h.q$ are injective on $W_i$.
\end{lemma}

\begin{proof}
Take $i$ so large that $K\cup L\cup KL\subset W_i$. Since $p\in M_i[W_i^4]$, the local multiplication identities at $p$ hold for every product used below. In particular,
\[
 h.(g.p)=(hg).p\qquad(\text{for every } h\in K).
\]
Put $q=g.p$. Next we check $q\in A_i=M_i[W_i^2]$.

For $a\in W_i^2$ we have $ag\in W_i^3\subset W_i^4$, and so
\[
 a.q=a.(g.p)=(ag).p.
\]
Right multiplication by $g$ is a homeomorphism from $W_i^2$ onto $W_i^2g$, while the orbit map $u\mapsto u.p$ is a homeomorphism on $W_i^4$. Their composition then shows that $a\mapsto a.q$ is a homeomorphism from $W_i^2$ onto an open neighborhood of $q$. If $a,b,ab\in W_i^2$, then the same calculation, now using $b g,abg\in W_i^4$, gives
\[
 a.(b.q)=a.(bg.p)=(abg).p=(ab).(g.p)=(ab).q.
\]
So $q\in M_i[W_i^2]=A_i$. The orbit maps at $p$ and $q$ are also injective on $W_i\subset W_i^2$.
\end{proof}

\section{The Poisson configuration space and completed root views}



\begin{definition}[Poisson point process]
	Let $(X,\mathcal B,\lambda)$ be a $\sigma$-finite measure space.
	A Poisson point process with intensity measure $\lambda$ is a random
	locally finite counting measure $\eta$ on $X$ such that for every
	finite collection of pairwise disjoint measurable sets
	$A_1,\dots,A_k$ with $\lambda(A_j)<\infty$,
	\[
	\eta(A_1),\dots,\eta(A_k)
	\]
	are independent random variables and
	\[
	\eta(A_j)\sim {\rm Pois}(\lambda(A_j)).
	\]
	We write $\PPP(\lambda)$ for this law.
\end{definition}
Let $\cN(G)$ be the space of locally finite integer-valued Radon measures on $G$, with the vague topology and its Borel $\sigma$-algebra. An element can be written as
\[
\omega = \sum_{x \in P} \delta_x,
\]
where $P$ is a locally finite point set allowing multiplicities. The vague topology is characterized by
\[
\omega_n \to \omega \quad \Leftrightarrow \quad \int f \, d\omega_n \to \int f \, d\omega, \quad \forall f \in C_c(G).
\]
Equipped with this topology, $\cN(G)$ is a Polish space. For $g\in G$ define
\begin{equation}\label{eq:right-shift}
 (g\omega)(E):=\omega(Eg),\qquad E\subset G\text{ Borel}.
\end{equation}
This is a continuous left action. Unimodularity implies that $\mu_\alpha:=\PPP(\alpha\Haar)$ is invariant.

We fix a bounded compatible metric $\rho_X$ on $\cN(G)$ with the following locality property:

\begin{equation}\label{eq:local-metric}
 \text{for every }\eta>0\text{ there is compact }K\subset G\text{ such that }
 \omega|_K=\omega'|_K\Longrightarrow \rho_X(\omega,\omega')<\eta.
\end{equation}

	To define such a metric, choose a compact exhaustion
\[
K_1\subset\operatorname{int}K_2\subset\cdots,
\qquad \bigcup_{m\geq1}K_m=G.
\]
For every $m$, choose a cutoff $\chi_m\in C_c(G)$ satisfying
$0\leq\chi_m\leq1$, $\chi_m\equiv1$ on $K_m$, and
$\operatorname{supp}(\chi_m)\subseteq K_{m+1}$.  Also choose a countable
family dense, in the uniform norm, in the continuous functions supported in
$K_m$. Enumerate the union of all these dense families and all the cutoffs as
$(f_n)_{n\geq1}$, arranging that every finite initial segment has supports in
one compact set.

This family determines convergence. In fact, if
$\int f_n\,d\omega_j\to\int f_n\,d\omega$ for every $n$, then for
each $m$ the convergence of the cutoff coordinate gives
\[
\sup_j\omega_j(K_m)
\leq \sup_j\int\chi_m\,d\omega_j<\infty.
\]
The resulting uniform local-mass bound allows uniform approximation on
$K_m$ to pass from the chosen dense family to every $f\in C_c(G)$.
So we have
\[
\rho_X(\omega,\omega')
:=\sum_{n\geq1}2^{-n}
\frac{\abs{\int f_n\dd\omega-\int f_n\dd\omega'}}
{1+\abs{\int f_n\dd\omega-\int f_n\dd\omega'}}
\]
is a bounded compatible metric on $\cN(G)$. It has the required locality
property: given $\eta>0$, choose $N$ so large that
$\sum_{n>N}2^{-n}<\eta$ and let $K$ contain the supports of
$f_1,\dots,f_N$. If $\omega|_K=\omega'|_K$, then the first $N$ summands
vanish and $\rho_X(\omega,\omega')<\eta$.

For $p\in A_i$, let
\[
 \theta_{i,p}:W_i\longrightarrow W_i.p,\qquad h\longmapsto h.p.
\]
This is a homeomorphism. Its inverse gives the coordinate of a point relative to the root $p$.

\begin{lemma}[Borel root-coordinate map]\label{lem:borel-root}
For each $i$, the set
\[
 \mathscr R_i:=\{(p,x)\in A_i\times M_i:x\in W_i.p\}
\]
is Borel, and the map
\[
 c_i:\mathscr R_i\longrightarrow W_i,
 \qquad c_i(p,x)=\theta_{i,p}^{-1}(x),
\]
is Borel. It also follows that if $\cN(M_i)$ denotes the standard Borel space of locally finite integer-valued Radon measures on $M_i$, then
\[
 (p,\xi)\longmapsto (c_i(p,\cdot))_*(\xi|_{W_i.p})
\]
is a Borel map from $A_i\times\cN(M_i)$ to $\cN(W_i)$.
\end{lemma}

\begin{proof}
The set $A_i=M_i[W_i^2]$ is Borel by the standard measurability of good sets in a local $G$-space. Consider
\[
 T_i:W_i\times A_i\longrightarrow A_i\times M_i,
 \qquad T_i(h,p)=(p,h.p).
\]
The partial action is continuous on its domain, so $T_i$ is continuous. It is injective because, for every $p\in A_i=M_i[W_i^2]$, the orbit map $h\mapsto h.p$ is injective on $W_i$. Both $W_i\times A_i$ and $A_i\times M_i$ are standard Borel spaces. The Lusin--Souslin theorem then implies that the image
\[
 \mathscr R_i=T_i(W_i\times A_i)
 =\{(p,x):p\in A_i,\ x\in W_i.p\}
\]
is Borel and that $T_i^{-1}:\mathscr R_i\to W_i\times A_i$ is Borel. The first coordinate of $T_i^{-1}(p,x)$ is exactly $c_i(p,x)$, so $c_i$ is Borel.

For the claim concerning point measures, we use the following elementary fact about integration against a varying point measure. If $P$ and $X$ are standard Borel spaces and $u:P\times X\to[0,\infty]$ is Borel, then
\begin{equation}\label{eq:point-measure-param-int}
 (p,\xi)\longmapsto\int_Xu(p,x)\,d\xi(x)
\end{equation}
is Borel on $P\times\cN(X)$. To verify this, first take $u=\1_{A\times B}$. Then \eqref{eq:point-measure-param-int} equals $\1_A(p)\xi(B)$, which is Borel because the evaluation maps $\xi\mapsto\xi(B)$ generate the standard Borel structure on $\cN(X)$. The class of nonnegative Borel functions for which the claim holds is closed under nonnegative linear combinations and monotone limits. The monotone class theorem then proves the claim for every nonnegative Borel $u$. Bounded signed functions follow by applying the result to the positive and negative parts.

Fix $f\in C_c(W_i)$. Extend the function
\[
 (p,x)\longmapsto \1_{\mathscr R_i}(p,x)f(c_i(p,x))
\]
by zero outside $\mathscr R_i$. It is a bounded Borel function on $A_i\times M_i$. For every $(p,\xi)$,
\begin{align*}
 &\int_{W_i}f(h)\,d\bigl[(c_i(p,\cdot))_*(\xi|_{W_i.p})\bigr](h)\\
 &\qquad=\int_{M_i}\1_{\mathscr R_i}(p,x)f(c_i(p,x))\,d\xi(x).
\end{align*}
The right-hand side is Borel in $(p,\xi)$ by the previous integration fact. It is finite because $\operatorname{supp}f$ is compact and its image $(\operatorname{supp}f).p$ is compact, while $\xi$ is locally finite.

Finally, choose a countable convergence-determining family $(f_m)$ in $C_c(W_i)$. The standard Borel structure of $\cN(W_i)$ is generated by the maps
\[
 \eta\longmapsto\int f_m\,d\eta.
\]
Since all these coordinate maps are Borel after composition with
$(p,\xi)\mapsto(c_i(p,\cdot))_*(\xi|_{W_i.p})$, the latter map is Borel as claimed.
\end{proof}

Let
\[
 \Xi_i\sim\PPP(\alpha\vol_{M_i})
\]
be a Poisson point process on $M_i$, and independently let
\[
 Z_i\sim\PPP(\alpha\Haar)
\]
be a Poisson point process on $G$. Define the random map
\[
 \Phi_i:M_i\longrightarrow\cN(G)
\]
by
\begin{equation}\label{eq:completed-root}
 \Phi_i(p):=
 \begin{cases}
 (\theta_{i,p}^{-1})_*(\Xi_i|_{W_i.p})+Z_i|_{G\setminus W_i},&p\in A_i,\\[2mm]
 Z_i,&p\notin A_i.
 \end{cases}
\end{equation}
By \cref{lem:borel-root}, the pulled-back interior point measure is jointly Borel on $A_i\times\cN(M_i)$. Restriction of a point measure to a fixed Borel set, pushforward under a Borel map, and superposition of point measures are Borel operations on the corresponding point-measure spaces. Since $A_i$ is Borel, the two cases in \eqref{eq:completed-root} then combine to show that the map from the underlying Poisson probability space times $M_i$ into $\cN(G)$, $(\omega,p)\mapsto\Phi_i(\omega,p)$, is jointly Borel.

\begin{proposition}[Exact Poisson completion lemma]\label{prop:exact-marginal}
For every $i$ and every $p\in M_i$,
\[
 \Phi_i(p)\sim\mu_\alpha.
\]
In particular, for every bounded Borel $F:\cN(G)\to\mathbb R$,
\begin{equation}\label{eq:exact-expectation}
 \mathbb E[F(\Phi_i(p))]=\int F\dd\mu_\alpha.
\end{equation}
\end{proposition}

\begin{proof}
If $p\notin A_i$, the claim is clear. Suppose $p\in A_i$. The restriction $\Xi_i|_{W_i.p}$ is a Poisson process with intensity $\alpha\vol_{M_i}|_{W_i.p}$. By the defining property of the canonical measure and the mapping theorem for Poisson processes, its pullback through $\theta_{i,p}$ is a Poisson process on $W_i$ with intensity $\alpha\Haar|_{W_i}$. It is independent of $Z_i|_{G\setminus W_i}$, which is Poisson with intensity $\alpha\Haar|_{G\setminus W_i}$. The superposition on the disjoint partition $G=W_i\sqcup(G\setminus W_i)$ is then $\PPP(\alpha\Haar)$.

The exterior term is essential for this exact identity. Without it, the
truncated view would be supported on configurations with no points outside
$W_i$, a $\mu_\alpha$-null set. An arbitrary measurable factor map is not
controlled on such a null set, while the completed view has the correct law
on the whole configuration space.
\end{proof}

\begin{proposition}[Exact local covariance on the deep core]\label{prop:local-covariance}
Let $K,L\subset G$ be compact. For all sufficiently large $i$, for every $p\in D_i$ and $g\in L$,
\begin{equation}\label{eq:local-covariance}
 \Phi_i(g.p)|_K=(g\Phi_i(p))|_K
\end{equation}
almost surely.
\end{proposition}

\begin{proof}
For large $i$, \cref{lem:change-root} applies, $K\cup KL\subset W_i$, and $q=g.p$ belongs to $A_i$. The exterior process $Z_i|_{G\setminus W_i}$ contributes no points to $K$ and no points to $Kg$. So both sides of \eqref{eq:local-covariance} are determined by $\Xi_i$.

For a Borel set $E\subset K$, the left-hand side counts $x\in\Xi_i$ for which $x=h.q$ with $h\in E$. By \cref{lem:change-root}, $h.q=(hg).p$. So this equals the number of root-$p$ coordinates in $Eg$, that is
\[
 \Phi_i(q)(E)=\Phi_i(p)(Eg)=(g\Phi_i(p))(E)
\]
by \eqref{eq:right-shift}. So the restrictions agree, as claimed.
\end{proof}

\section{Exponential concentration for local observables}

Let $K \subseteq G$ be a compact set. A bounded Borel function
\[
F : \cN(G) \to \mathbb{R}
\]
is called $K$-local if $F(\omega)$ depends only on $\omega|_K$, for any $\omega\in \cN(G)$. That is, if
$\omega|_K = \omega'|_K$,
then
$F(\omega) = F(\omega')$.


We begin with the geometric overlap estimate.

\begin{lemma}\label{lem:overlap}
Let $K \subseteq G$ be a compact set. For all sufficiently large $i$ and every $z\in M_i$,
\begin{equation}\label{eq:overlap}
 \vol_{M_i}\bigl(\{p\in D_i:z\in K.p\}\bigr)\leq \Haar(K).
\end{equation}
\end{lemma}

\begin{proof}
Let $E_z:=\{p\in D_i:z\in K.p\}$. If $E_z=\varnothing$ there is nothing to prove. Otherwise choose $p_0\in E_z$ and $k_0\in K$ with $z=k_0.p_0$. By \cref{lem:change-root}, applied with the fixed compact set $K$, we have $z\in A_i=M_i[W_i^2]$ for all sufficiently large $i$.

For $p\in E_z$, injectivity of the orbit map at $p$ gives a unique $k(p)\in K$ such that $z=k(p).p$. The inverse axiom for the partial action gives
\[
 p=k(p)^{-1}.z.
\]
So we have
\[
 E_z\subset K^{-1}.z.
\]
Since $z\in M_i[W_i^2]$ and $K^{-1}\subset W_i^2$ for large $i$, the single orbit chart centered at $z$ is defined and injective on all of $K^{-1}$. The canonical-measure formula in this chart gives
\[
 \vol_{M_i}(E_z)\leq \vol_{M_i}(K^{-1}.z)
 =\Haar(K^{-1})=\Haar(K),
\]
where the last equality follows from unimodularity.
\end{proof}

We use the following consequence of Wu's modified logarithmic Sobolev inequality \cite{Wu}. Since the precise Bernstein form is needed later, we include the complete Herbst argument; the logarithmic Sobolev inequality used in the first step is stated in \cite[Theorem~1.1]{BachmannPeccati}.

\begin{lemma}[Poisson bounded differences]\label{lem:poisson-bd}
Let $\eta$ be a Poisson process on a finite measure space $(S,\lambda)$, and let $H=H(\eta)$ be bounded. Suppose that
\[
 \abs{H(\xi+\delta_z)-H(\xi)}\leq c
\]
for every locally finite integer-valued $\xi$ and every $z\in S$. Then, for every $t>0$,
\begin{equation}\label{eq:poisson-bd}
 \mathbb P\bigl(\abs{H-\mathbb EH}\geq t\bigr)
 \leq2\exp\left(-\frac{t^2}{2\lambda(S)c^2+2ct/3}\right).
\end{equation}
\end{lemma}

\begin{proof}
Write
\[
 D_zH(\xi):=H(\xi+\delta_z)-H(\xi).
\]
If $c=0$, then $H(\xi+\delta_z)=H(\xi)$ for every $\xi$ and $z$. Since $\lambda(S)<\infty$, the Poisson process has finitely many points almost surely, and repeated removal of its points shows that $H(\eta)=H(0)$ almost surely. The desired estimate is then trivial. From now on, assume $c>0$. We first prove the upper-tail estimate. For $\theta\geq0$, put
\[
 L(\theta):=\log\mathbb E e^{\theta H}.
\]
Because $H$ is bounded, $L$ is differentiable and all exponential moments below are finite. Wu's modified logarithmic Sobolev inequality, in the form of \cite[Theorem~1.1]{BachmannPeccati}, states that
\begin{equation}\label{eq:wu-mlsi-applied}
 \operatorname{Ent}(e^{\theta H})
 \leq
 \mathbb E\left[e^{\theta H}
 \int_S\psi(\theta D_zH)\,d\lambda(z)\right],
\end{equation}
where $\operatorname{Ent}(X):=\mathbb E[X\log X]-\mathbb E[X]\log\mathbb E(X)$ and $\psi(u):=ue^u-e^u+1$.
The function $\psi$ is nonnegative, and for every $a\geq0$ one has
\[
 \max_{|u|\leq a}\psi(u)=\psi(a).
\]
In fact, $\psi'(u)=ue^u$, so $\psi$ decreases on $(-\infty,0]$ and increases on $[0,\infty)$. Also,
\[
 \psi(a)-\psi(-a)=2(a\cosh a-\sinh a)\geq0,
\]
because the derivative of $a\cosh a-\sinh a$ is $a\sinh a\geq0$ and its value at $a=0$ is zero. Since $|D_zH|\leq c$, \eqref{eq:wu-mlsi-applied} gives
\begin{equation}\label{eq:entropy-bound-poisson}
 \operatorname{Ent}(e^{\theta H})
 \leq \lambda(S)\psi(\theta c)\mathbb E e^{\theta H}.
\end{equation}
But we also have
\[
 \frac{\operatorname{Ent}(e^{\theta H})}{\mathbb E e^{\theta H}}
 =\theta L'(\theta)-L(\theta).
\]
So we have
\begin{equation}\label{eq:Herbst-differential}
 \frac{d}{d\theta}\left(\frac{L(\theta)}{\theta}\right)
 =\frac{\theta L'(\theta)-L(\theta)}{\theta^2}
 \leq\lambda(S)\frac{\psi(\theta c)}{\theta^2}
 \qquad(\theta>0).
\end{equation}
Since $L(\theta)/\theta\to\mathbb EH$ as $\theta\downarrow0$, integration of \eqref{eq:Herbst-differential} from $0$ to $\theta$ gives
\begin{equation}\label{eq:poisson-mgf-bound}
 \log\mathbb E e^{\theta(H-\mathbb EH)}
 \leq\lambda(S)\bigl(e^{\theta c}-1-\theta c\bigr).
\end{equation}
Here we used the identity
\[
 \frac{d}{d\theta}
 \left(\frac{e^{\theta c}-1-\theta c}{\theta}\right)
 =\frac{\psi(\theta c)}{\theta^2}.
\]

For $u>0$, Chernoff's inequality and \eqref{eq:poisson-mgf-bound} imply
\[
 \mathbb P(H-\mathbb EH\geq u)
 \leq\inf_{\theta>0}
 \exp\left[-\theta u+\lambda(S)(e^{\theta c}-1-\theta c)\right].
\]
If $\lambda(S)=0$, the random process is almost surely empty and the desired estimate is trivial. Assume $\lambda(S)>0$ and set
\[
 v:=\frac{u}{\lambda(S)c}.
\]
The minimizing value is $\theta=c^{-1}\log(1+v)$, and substitution gives the Bennett bound
\begin{equation}\label{eq:Bennett-upper}
 \mathbb P(H-\mathbb EH\geq u)
 \leq\exp[-\lambda(S)h(v)],
 \qquad
 h(v):=(1+v)\log(1+v)-v.
\end{equation}
The elementary inequality
\begin{equation}\label{eq:h-Bernstein}
 h(v)\geq\frac{v^2}{2(1+v/3)},\qquad v\geq0,
\end{equation}
then gives
\begin{equation}\label{eq:poisson-upper-Bernstein}
 \mathbb P(H-\mathbb EH\geq u)
 \leq
 \exp\left(-\frac{u^2}{2\lambda(S)c^2+2cu/3}\right).
\end{equation}
Apply the same argument to $-H$. Its add-one differences satisfy
\[
 |D_z(-H)|=|D_zH|\leq c,
\]
so
\[
 \mathbb P(\mathbb EH-H\geq u)
 \leq
 \exp\left(-\frac{u^2}{2\lambda(S)c^2+2cu/3}\right).
\]
Adding the two one-sided bounds proves \eqref{eq:poisson-bd}.
\end{proof}

\begin{proposition}[Concentration of local root averages]\label{prop:local-concentration}
Let $K\subset G$ be compact and let $F:\cN(G)\to[-1,1]$ be $K$-local. Define
\[
 \mathcal A_i(F):=\frac1{V_i}\int_{D_i}F(\Phi_i(p))\dd\vol_{M_i}(p).
\]
Then
\begin{equation}\label{eq:local-mean}
 \mathbb E\mathcal A_i(F)=d_i\int F\dd\mu_\alpha.
\end{equation}
Also, for every $t>0$ there is $c=c(t,K,\alpha)>0$ such that, for all sufficiently large $i$,
\begin{equation}\label{eq:local-concentration}
 \mathbb P\left(\abs{\mathcal A_i(F)-\mathbb E\mathcal A_i(F)}>t\right)
 \leq2e^{-cV_i}.
\end{equation}
The same claim holds for a function taking values in $[-1,1]$ that is
$K$-local in each coordinate of two independent completed root-view models.
\end{proposition}

\begin{proof}
The expectation formula follows from \cref{prop:exact-marginal} and Fubini's theorem.
\begin{align*}
	\mathbb E\mathcal A_i(F) &= \frac{1}{V_i} \int_{D_i} \mathbb E[F(\Phi_i(p))] \dd\vol_{M_i}(p) \\
	&= \frac{1}{V_i} \int_{D_i} \left( \int F \dd\mu_\alpha \right) \mathrm{d}\vol_{M_i}(p) \\
	&= \frac{\vol_{M_i}(D_i)}{V_i} \int F \dd\mu_\alpha \\
	&= d_i \int F \dd\mu_\alpha.
\end{align*}

For large $i$, $K\subset W_i$. On $D_i$, the value $F(\Phi_i(p))$ depends only on the restriction of $\Xi_i$ to $K.p$ and is independent of $Z_i$.

Add a point at $z\in M_i$ to $\Xi_i$. The integrand can change only for roots $p\in D_i$ satisfying $z\in K.p$, and it changes by at most $2$. By \cref{lem:overlap},
\[
 \abs{D_z\mathcal A_i(F)}\leq\frac{2\vol_{M_i}\bigl(\{p\in D_i:z\in K.p\}\bigr)}{V_i}\leq\frac{2\Haar(K)}{V_i}=:c_i.
\]
The intensity measure of $\Xi_i$ has total mass $\alpha V_i$. Applying \cref{lem:poisson-bd} with $c=c_i$ and $\lambda(S)=\alpha V_i$ gives
\[
 \mathbb P\left(\abs{\mathcal A_i(F)-\mathbb E\mathcal A_i(F)}>t\right)
 \leq2\exp\left(
 -\frac{t^2V_i}{8\alpha\Haar(K)^2+4t\Haar(K)/3}
 \right).
\]
This gives the stated exponential bound. For a bounded function that is $K$-local in each coordinate of two independent models, view the pair $(\Xi_i,\Xi_i')$ as one Poisson process on the disjoint union $M_i\sqcup M_i$. Its intensity has total mass $2\alpha V_i$. Adding one point in either copy changes the normalized root average only at roots $p$ whose corresponding window $K.p$ contains that point, so the same overlap argument gives the add-one bound $2\Haar(K)/V_i$. Applying \cref{lem:poisson-bd} gives explicitly
\[
 2\exp\left(
 -\frac{t^2V_i}{16\alpha\Haar(K)^2+4t\Haar(K)/3}
 \right),
\]
which is $2e^{-cV_i}$ for a constant $c=c(t,K,\alpha)>0$.
\end{proof}

\begin{remark}[Uniformity of the concentration bound]
The exponential rate in \cref{prop:local-concentration} is uniform over all
$K$-local observables $F$ with $\|F\|_\infty\leq1$. In particular, it depends
only on the window $K$, the intensity $\alpha$, and the deviation parameter
$t$, and not on the particular observable. The analogous two-copy estimate
has the same uniformity.
\end{remark}

\section{Random microstates for arbitrary measurable factors}

Let
\[
 \pi:(\cN(G),\mu_\alpha)\longrightarrow(Y,\nu)
\]
be a measurable factor, where $(Y,\nu)$ is a compact metrizable continuous model as given by \cref{prop:compact-model}. Fix a compatible metric $\rho_Y\leq1$.

Choose a Borel representative of the measurable factor map $\pi$. For every fixed $g\in G$, the factor property gives
\[
 \pi(g\omega)=g\pi(\omega)
\]
for $\mu_\alpha$-almost every $\omega$. The failure set
\[
 N:=\{(g,\omega)\in G\times\cN(G):
 \pi(g\omega)\neq g\pi(\omega)\}
\]
is Borel because both actions and $\pi$ are Borel. Every vertical section $N_g$ has $\mu_\alpha$-measure zero. Since Haar measure is $\sigma$-finite, Tonelli's theorem gives
\[
 (\Haar\times\mu_\alpha)(N)=0.
\]
So our Borel representative satisfies
\begin{equation}\label{eq:factor-ae}
 \pi(g\omega)=g\pi(\omega)
\end{equation}
for $(\Haar\times\mu_\alpha)$-almost every $(g,\omega)$. This joint null-set statement is the form needed for the averaged equivariance estimates below. Define
\[
 \Psi_i:=\pi\circ\Phi_i:M_i\longrightarrow Y.
\]
\begin{proposition}
\label{prop:factor-marginal}
For every $i$ and every $p\in M_i$,
\[
 \Psi_i(p)\sim\nu.
\]
In other words, for every bounded Borel function $f:Y\to\mathbb R$,
\[
 \mathbb E[f(\Psi_i(p))]=\int_Y f\,d\nu.
\]
\end{proposition}

\begin{proof}
By \cref{prop:exact-marginal}, $\Phi_i(p)$ has law $\mu_\alpha$. This gives
\[
 (\pi\circ\Phi_i(p))_*\mathbb P=\pi_*\mu_\alpha=\nu.
\]
\end{proof}

\subsection{Empirical convergence}


\begin{lemma}\label{lem:local-L1}
For every bounded Borel function $H:\cN(G)\to\mathbb R$ and every $\eta>0$, there are a compact $K\subset G$ and a bounded $K$-local Borel function $H_K$ such that
\[
 \norm{H-H_K}_{L^1(\mu_\alpha)}<\eta,
 \qquad \norm{H_K}_\infty\leq\norm{H}_\infty.
\]
\end{lemma}

\begin{proof}
Choose an increasing compact exhaustion
\[
 K_1\subseteq K_2\subseteq\cdots,
 \qquad \bigcup_{n\geq1}K_n=G,
\]
and let $\mathscr F_n$ be the $\sigma$-algebra generated by the restriction map
\[
 r_{K_n}:\omega\longmapsto\omega|_{K_n}.
\]
Compact restrictions generate the Borel $\sigma$-algebra of the configuration
space, so
\[
 \bigvee_{n\geq1}\mathscr F_n=\cB(\cN(G)).
\]
Let $H_n:=\mathbb E(H\mid\mathscr F_n)$.
The martingale convergence theorem gives $H_n\to H$ in $L^1(\mu_\alpha)$. Conditional expectation is a contraction on $L^\infty$, so
\[
 \|H_n\|_\infty\leq\|H\|_\infty.
\]
Choose $n$ such that $\|H-H_n\|_1<\eta$. Since $\mathscr F_n$ is generated by the Borel restriction map
\[
 r_{K_n}:\cN(G)\to\cN(K_n),\qquad\omega\mapsto\omega|_{K_n},
\]
the Doob--Dynkin lemma supplies a Borel function $\widetilde H_n$ on $\cN(K_n)$ such that
\[
 H_n(\omega)=\widetilde H_n(\omega|_{K_n})
\]
for $\mu_\alpha$-almost every $\omega$. After truncating $\widetilde H_n$ to the interval $[-\|H\|_\infty,\|H\|_\infty]$, this equality and the norm bound remain valid. Define
\[
 H_K(\omega):=\widetilde H_n(\omega|_{K_n}),\qquad K:=K_n.
\]
Then $H_K$ is Borel, $K$-local, satisfies $\|H_K\|_\infty\leq\|H\|_\infty$, and obeys $\|H-H_K\|_1<\eta$.
\end{proof}

\begin{proposition}\label{prop:empirical-factor}
For every weak-star neighborhood $\cO$ of $\nu$, when $i \to \infty$,
\[
 \mathbb P\bigl((\Psi_i)_*\dbar_i\in\cO\bigr)\longrightarrow1.
\]
\end{proposition}

\begin{proof}
It is enough to treat a basic weak-star neighborhood
\[
 \cO=\left\{\lambda\in\Prob(Y):
 \left|\int f_j\,d\lambda-\int f_j\,d\nu\right|<\tau
 \text{ for }1\leq j\leq m\right\},
\]
where $f_j\in C(Y)$ and $\|f_j\|_\infty\leq1$. Put $H_j:=f_j\circ\pi$. To prove convergence in probability, fix an
arbitrary $\xi>0$. Choose $\eta>0$ so small that
\begin{equation}\label{eq:empirical-eta-choice}
 \sqrt\eta+2\eta<\tau/3,
 \qquad m\sqrt\eta<\xi/2.
\end{equation}
By \cref{lem:local-L1}, for each $j$ there is a compact window and a local approximation with $L^1$ error less than $\eta$. Replacing the finitely many windows by their union, we get one compact $K\subset G$ and $K$-local Borel functions $H_{j,K}$ satisfying
\[
 \|H_j-H_{j,K}\|_{L^1(\mu_\alpha)}<\eta,
 \qquad \|H_{j,K}\|_\infty\leq1
 \quad(1\leq j\leq m).
\]

Define
\[
 E_{i,j}:=\frac1{V_i}\int_{M_i}
 |H_j(\Phi_i(p))-H_{j,K}(\Phi_i(p))|\,d\vol_{M_i}(p).
\]
The exact marginal identity and Tonelli's theorem give
\[
 \mathbb E E_{i,j}
 =\int_{\cN(G)}|H_j-H_{j,K}|\,d\mu_\alpha<\eta.
\]
So we have
\begin{equation}\label{eq:empirical-approx-prob}
 \mathbb P(E_{i,j}>\sqrt\eta)<\sqrt\eta
\end{equation}
by Markov's inequality.

Set
\[
 J_{i,j}:=\frac1{V_i}\int_{D_i}H_{j,K}(\Phi_i(p))\,d\vol_{M_i}(p),
 \qquad m_{j,K}:=\int H_{j,K}\,d\mu_\alpha.
\]
By \cref{prop:local-concentration},
\[
 J_{i,j}-d_i m_{j,K}\xrightarrow{\mathbb P}0.
\]
Since $d_i\to1$, we get $J_{i,j}\to m_{j,K}$ in probability.
So there is $i_1$ such that, for every $i\geq i_1$,
\begin{equation}\label{eq:empirical-local-prob}
 \sum_{j=1}^m
 \mathbb P\left(|J_{i,j}-m_{j,K}|>\tau/3\right)<\xi/2.
\end{equation}
The contribution from the complement of the deep core satisfies
\[
 \left|\frac1{V_i}\int_{M_i\setminus D_i}
 H_{j,K}(\Phi_i(p))\,d\vol_{M_i}(p)\right|
 \leq1-d_i.
\]
Also,
\[
 |m_{j,K}-\int H_j\,d\mu_\alpha|<\eta.
\]
On the event $E_{i,j}\leq\sqrt\eta$,
\begin{align*}
 &\left|\frac1{V_i}\int_{M_i}H_j(\Phi_i(p))\,d\vol_{M_i}(p)
 -\int H_j\,d\mu_\alpha\right|\\
 &\quad\leq E_{i,j}+|J_{i,j}-m_{j,K}|+(1-d_i)+\eta.
\end{align*}
Increase $i_1$ if necessary so that
$1-d_i<\tau/3-\sqrt\eta-\eta$ for every $i\geq i_1$, which is possible by
\eqref{eq:empirical-eta-choice}. For such $i$, failure of any one of the $m$
required moment inequalities is contained in the union of the events
\[
 \{E_{i,j}>\sqrt\eta\}
 \quad\text{and}\quad
 \{|J_{i,j}-m_{j,K}|>\tau/3\},
 \qquad 1\leq j\leq m.
\]
By \eqref{eq:empirical-approx-prob},
\eqref{eq:empirical-local-prob}, and the union bound,
\[
 \mathbb P\bigl((\Psi_i)_*\dbar_i\notin\cO\bigr)
 \leq m\sqrt\eta+\xi/2<\xi.
\]
Here we used $\int H_j\,d\mu_\alpha=\int f_j\,d\nu$. Since $\xi>0$ was
arbitrary, the required probability tends to one.
\end{proof}

\subsection{Approximate equivariance}

Recall from \eqref{eq:defect} that a map $\psi:M_i\to Y$ is $(L,\delta)$-equivariant on average if
\[
 \frac1{\Haar(L)}\int_L\Delta_{\rho_Y,M_i}(g,\psi)\dd\Haar(g)<\delta.
\]

\begin{proposition}\label{prop:avg-equivariance}
For every precompact identity neighborhood $L\subset G$, when $i \to \infty$,
\[
 \frac1{\Haar(L)}\int_L
 \Delta_{\rho_Y,M_i}(g,\Psi_i)\dd\Haar(g)
 \xrightarrow{\mathbb P}0.
\]
\end{proposition}

\begin{proof}
Fix $\eta>0$. By Lusin's theorem, there is a compact set $C\subset\cN(G)$ such that
\[
 \mu_\alpha(C)>1-\eta
\]
and $\pi|_C$ is continuous. Since $C$ is compact, the restriction $\pi|_C$ is uniformly continuous. So there is $\delta_X>0$ such that
\begin{equation}\label{eq:lusin-modulus}
 x,x'\in C,\quad \rho_X(x,x')<\delta_X
 \quad\Longrightarrow\quad
 \rho_Y(\pi(x),\pi(x'))<\eta.
\end{equation}
By the locality property \eqref{eq:local-metric}, choose a compact set $K\subset G$ such that
\begin{equation}\label{eq:agreement-K-deltaX}
 x|_K=x'|_K\quad\Longrightarrow\quad\rho_X(x,x')<\delta_X.
\end{equation}

The set $L$ in the statement is precompact. Apply \cref{prop:local-covariance} to the two compact sets $K$ and $\overline L$. For all sufficiently large $i$, every $p\in D_i$ and every $g\in L$ satisfy
\[
 \Phi_i(g.p)|_K=(g\Phi_i(p))|_K.
\]
Combining this identity with \eqref{eq:agreement-K-deltaX}, we get
\begin{equation}\label{eq:root-close-lusin}
 \rho_X(\Phi_i(g.p),g\Phi_i(p))<\delta_X
 \qquad(p\in D_i,\ g\in L).
\end{equation}

We also justify carefully the use of the factor identity under the random root-view law. Let $N\subset G\times\cN(G)$ be the product-null failure set introduced above. For fixed $p\in M_i$, the exact marginal property gives
\begin{align*}
 &\int_L\mathbb P\bigl((g,\Phi_i(p))\in N\bigr)\,d\Haar(g)\\
 &\qquad=\int_L\int_{\cN(G)}\1_N(g,\omega)\,d\mu_\alpha(\omega)\,d\Haar(g)=0.
\end{align*}
Integrating once more over $p$ shows that, outside a
$(\Haar|_L\times\mathbb P\times\dbar_i)$-null set,
\begin{equation}\label{eq:factor-identity-random}
 g\Psi_i(p)=g\pi(\Phi_i(p))=\pi(g\Phi_i(p)).
\end{equation}

By Fubini's theorem, for $\Haar|_L\times\dbar_i$-almost every pair $(g,p)\in L\times D_i$, identity \eqref{eq:factor-identity-random} holds almost surely in the Poisson randomness. Fix such a pair $(g,p)$. If both configurations
\[
 \Phi_i(g.p)\quad\text{and}\quad g\Phi_i(p)
\]
belong to $C$, then \eqref{eq:root-close-lusin}, \eqref{eq:lusin-modulus}, and \eqref{eq:factor-identity-random} give
\begin{align*}
 \rho_Y(\Psi_i(g.p),g\Psi_i(p))
 =\rho_Y\bigl(\pi(\Phi_i(g.p)),\pi(g\Phi_i(p))\bigr)<\eta.
\end{align*}
For each fixed $p$ and $g$, the random variable $\Phi_i(g.p)$ has law $\mu_\alpha$ by \cref{prop:exact-marginal}. The variable $g\Phi_i(p)$ also has law $\mu_\alpha$, because $\Phi_i(p)\sim\mu_\alpha$ and $\mu_\alpha$ is $G$-invariant. They need not be independent, but the union bound is sufficient:
\[
 \mathbb P\bigl(\Phi_i(g.p)\notin C
 \text{ or }g\Phi_i(p)\notin C\bigr)
 \leq2\eta.
\]
Since $\rho_Y\leq1$, we get
\begin{equation}\label{eq:expected-core-defect}
 \mathbb E\,\rho_Y(\Psi_i(g.p),g\Psi_i(p))
 \leq \eta+2\eta=3\eta
 \qquad(p\in D_i,\ g\in L).
\end{equation}

For $p\in D_i=M_i[W_i^4]$ and large $i$, every $g\in L\subset W_i$ acts at $p$, so there is no undefined-domain penalty on $D_i$. On $M_i\setminus D_i$, the sum of the metric integrand on the defined part and the undefined-domain penalty is bounded by $1$. Using Tonelli's theorem and \eqref{eq:expected-core-defect}, we get
\begin{align*}
 &\mathbb E\left[
 \frac1{\Haar(L)}\int_L
 \Delta_{\rho_Y,M_i}(g,\Psi_i)\,d\Haar(g)
 \right]\\
 &\quad\leq
 \frac1{\Haar(L)}\int_L
 \left(\int_{D_i}3\eta\,d\dbar_i(p)
 +\dbar_i(M_i\setminus D_i)\right)d\Haar(g)\\
 &\quad=3\eta d_i+(1-d_i)
 \leq3\eta+(1-d_i).
\end{align*}
Letting $i\to\infty$ and using $d_i\to1$ gives
\[
 \limsup_{i\to\infty}
 \mathbb E\left[
 \frac1{\Haar(L)}\int_L
 \Delta_{\rho_Y,M_i}(g,\Psi_i)\,d\Haar(g)
 \right]
 \leq3\eta.
\]
Since $\eta>0$ was arbitrary, the expectations tend to zero. Markov's inequality then implies convergence to zero in probability, as required.
\end{proof}

Combining averaged equivariance with empirical convergence gives the random microstates used in the packing argument. The comparison theorem will later transfer the resulting lower bound to the standard model spaces.

\begin{corollary}\label{cor:good-microstates}
Let $L\subset G$ be a precompact identity neighborhood, $\delta>0$, and let $\cO$ be a weak-star neighborhood of $\nu$. Then when $i \to \infty$,
\[
 \mathbb P\left(
 \Psi_i\in\Mapavg(M_i,Y,\rho_Y:L,\delta,\cO)
 \right)\longrightarrow1.
\]
\end{corollary}

\begin{proof}
By \cref{prop:avg-equivariance}, the averaged equivariance condition holds with probability tending to one, and by \cref{prop:empirical-factor} the empirical measure lies in $\cO$ with probability tending to one.
\end{proof}

\section{Exponential separation in nontrivial factors}
Assume from now on that the factor is nontrivial. So $\nu$ is not a point mass.

\begin{lemma}\label{lem:separated-compacts}
There exist compact sets $C_0,C_1\subset Y$ and constants $p_0,p_1,\kappa>0$ such that
\[
 \nu(C_j)=p_j>0\quad(j=0,1),
 \qquad
 \inf\{\rho_Y(y_0,y_1):y_j\in C_j\}\geq\kappa.
\]
\end{lemma}

\begin{proof}
Because $\nu$ is not a point mass, its support contains two distinct points $y_0$ and $y_1$. Put
\[
 d:=\rho_Y(y_0,y_1)>0
\]
and choose $0<r<d/4$. The open balls $B_Y(y_j,r)$ have positive $\nu$-measure because $y_j\in\operatorname{supp}\nu$. By inner regularity of the Borel probability measure $\nu$ on the compact metric space $Y$, there are compact sets
\[
 C_j\subset B_Y(y_j,r)
\]
with $p_j:=\nu(C_j)>0$. If $z_j\in C_j$, then
\[
 \rho_Y(z_0,z_1)
 \geq\rho_Y(y_0,y_1)-\rho_Y(z_0,y_0)-\rho_Y(z_1,y_1)
 >d-2r>d/2.
\]
So the conclusion holds with $\kappa:=d/2$.
\end{proof}

Define
\[
 a:Y\to\{0,1,*\},\qquad
 a(y)=\begin{cases}0,&y\in C_0,\\1,&y\in C_1,\\ *,&\text{otherwise}.
 \end{cases}
\]

\begin{lemma}[Finite-window approximation of measurable factor labels]\label{lem:local-label}
For every $\eta>0$ there are a compact $K\subset G$ and a $K$-local Borel map
\[
 b:\cN(G)\to\{0,1,*\}
\]
such that
\begin{equation}\label{eq:label-error}
 \mu_\alpha\{b\neq a\circ\pi\}<\eta.
\end{equation}
Also, if $\eta<\min(p_0,p_1)/2$, then
\begin{equation}\label{eq:label-masses}
 r_j:=\mu_\alpha\{b=j\}\geq p_j-\eta>0,
 \qquad j=0,1.
\end{equation}
\end{lemma}

\begin{proof}
Let
\[
 A_0:=(a\circ\pi)^{-1}(0),\qquad
 A_1:=(a\circ\pi)^{-1}(1),\qquad
 A_*:=(a\circ\pi)^{-1}(*).
\]
These three Borel sets form a partition of $\cN(G)$. Choose an increasing compact exhaustion $(K_n)$ of $G$ and write $\mathscr F_n:=\mathscr F_{K_n}$. Since $\bigvee_n\mathscr F_n=\cB(\cN(G))$, the martingale convergence theorem gives, for each $j\in\{0,1,*\}$, when $n \to \infty$,
\[
 u_{j,n}:=\mathbb E(\1_{A_j}\mid\mathscr F_n)
 \longrightarrow\1_{A_j}
 \quad\text{in }L^1(\mu_\alpha).
\]
The conditional probabilities may be chosen Borel and $K_n$-local, and they satisfy
\[
 u_{0,n}+u_{1,n}+u_{*,n}=1
\]
almost everywhere. The exceptional set is $\mathscr F_n$-measurable. Redefine the triple there to be $(1,0,0)$. This changes none of the $L^1$ classes, preserves Borelness and $K_n$-locality, and makes the displayed identity hold at every configuration. Choose $n$ so large that
\begin{equation}\label{eq:partition-martingale-error}
 2\sum_{j\in\{0,1,*\}}
 \|u_{j,n}-\1_{A_j}\|_1<\eta.
\end{equation}
Define $b(\omega)$ to be an index at which the three numbers
$u_{0,n}(\omega),u_{1,n}(\omega),u_{*,n}(\omega)$ attain their maximum, using the fixed order $0<1<*$ to break ties. Then $b$ is Borel and $K_n$-local.

We estimate its error. If $\omega\in A_j$ but $b(\omega)\neq j$, then some other conditional probability is at least $u_{j,n}(\omega)$. Since the three probabilities sum to one, this forces $u_{j,n}(\omega)\leq1/2$. This gives
\[
 \1_{A_j\cap\{b\neq j\}}(\omega)
 \leq2\1_{A_j}(\omega)(1-u_{j,n}(\omega)).
\]
After integration and summation over the three labels,
\begin{align*}
 \mu_\alpha\{b\neq a\circ\pi\}
 &\leq2\sum_j\int_{A_j}(1-u_{j,n})\,d\mu_\alpha\\
 &\leq2\sum_j\|u_{j,n}-\1_{A_j}\|_1<\eta
\end{align*}
by \eqref{eq:partition-martingale-error}. This proves \eqref{eq:label-error} with $K:=K_n$.

For $j=0,1$, the symmetric difference estimate gives
\[
 \mu_\alpha\{b=j\}
 \geq\mu_\alpha(A_j)-\mu_\alpha\{b\neq a\circ\pi\}
 >p_j-\eta.
\]
If $\eta<\min(p_0,p_1)/2$, these quantities are positive, proving \eqref{eq:label-masses}.
\end{proof}

\begin{remark}
This is the measurable analogue of finite-coordinate approximation for
Bernoulli factors. No continuity of the factor map is assumed: the argument
uses only that compact restrictions generate the standard Borel structure of
the Poisson configuration space.
\end{remark}

Let
$
 q_*:=\frac{p_0p_1}{4}>0.
$
Choose $\eta>0$ so small that
\begin{equation}\label{eq:eta-choice}
 \eta<\frac12\min(p_0,p_1),
 \qquad
 2\sqrt\eta<\frac{q_*}{2}.
\end{equation}
Apply \cref{lem:local-label} with this $\eta$, and fix the resulting compact window $K$ and local label $b$. Then \eqref{eq:label-masses} gives $r_j\geq p_j/2$ and so $r_0r_1\geq q_*$.

Let $(\Phi_i',\Psi_i')$ be an independent copy of $(\Phi_i,\Psi_i)$, built from independent Poisson processes $\Xi_i',Z_i'$. Define the local opposite-label statistic
\begin{equation}\label{eq:Zi}
 Z_i^\mathrm{opp}:=\frac1{V_i}\int_{D_i}
 \1_{\{(b(\Phi_i(p)),b(\Phi_i'(p)))\in\{(0,1),(1,0)\}\}}
 \dd\vol(p).
\end{equation}

\begin{proposition}\label{prop:opposite-label}
For all sufficiently large $i$,
\[
 \mathbb E Z_i^\mathrm{opp}=2d_i r_0r_1\geq2d_i q_*.
\]
There exists $\gamma=\gamma(K,\alpha,q_*)>0$ such that
\begin{equation}\label{eq:opposite-tail}
 \mathbb P(Z_i^\mathrm{opp}<q_*)\leq2e^{-\gamma V_i}
\end{equation}
for all sufficiently large $i$.
\end{proposition}

\begin{proof}
Fix $i$ large enough that $K\subset W_i$. For $p\in D_i$, the $K$-locality of $b$ and the definition of the completed root view imply that $b(\Phi_i(p))$ depends only on the restriction $\Xi_i|_{K.p}$; the exterior process contributes no point to the window $K$. The exact marginal property gives
\[
 b(\Phi_i(p))\sim b_*\mu_\alpha.
\]
The primed and unprimed constructions are independent, so the two labels at a fixed root are independent. This gives
\begin{align*}
 &\mathbb P\bigl((b(\Phi_i(p)),b(\Phi_i'(p)))
 \in\{(0,1),(1,0)\}\bigr)\\
 &\qquad=r_0r_1+r_1r_0=2r_0r_1.
\end{align*}
Tonelli's theorem now gives
\[
 \mathbb E Z_i^{\mathrm{opp}}
 =\frac{\vol_{M_i}(D_i)}{V_i}\,2r_0r_1
 =2d_i r_0r_1\geq2d_i q_*.
\]

Define on two configurations
\[
 F(\omega,\omega')
 :=\1_{\{(b(\omega),b(\omega'))\in\{(0,1),(1,0)\}\}}.
\]
This function takes values in $[0,1]$ and is local in the window $K$ in each coordinate. The two-model part of \cref{prop:local-concentration} then gives, for the fixed number $t=q_*/2$, a constant $\gamma>0$ such that
\[
 \mathbb P\left(
 |Z_i^{\mathrm{opp}}-\mathbb E Z_i^{\mathrm{opp}}|>q_*/2
 \right)\leq2e^{-\gamma V_i}
\]
for all sufficiently large $i$. Since $d_i\to1$, we may also assume $d_i\geq3/4$, and then
\[
 \mathbb E Z_i^{\mathrm{opp}}\geq2d_i q_*\geq3q_*/2.
\]
So the event $Z_i^{\mathrm{opp}}<q_*$ implies
\[
 |Z_i^{\mathrm{opp}}-\mathbb E Z_i^{\mathrm{opp}}|>q_*/2,
\]
which proves \eqref{eq:opposite-tail}.
\end{proof}

Define
\[
 R_i:=\frac1{V_i}\int_{M_i}
 \1_{\{b(\Phi_i(p))\neq a(\Psi_i(p))\}}
 \dd\vol(p),
\]
and define $R_i'$ analogously.

\begin{lemma}[Factor-label error]\label{lem:Ri}
One has
\[
 \mathbb E R_i<\eta,
 \qquad
 \mathbb P(R_i>\sqrt\eta)<\sqrt\eta.
\]
The same statements hold for $R_i'$.
\end{lemma}

\begin{proof}
For every fixed $p\in M_i$, the exact marginal property and the definition $\Psi_i(p)=\pi(\Phi_i(p))$ give
\begin{align*}
 &\mathbb P\bigl(b(\Phi_i(p))\neq a(\Psi_i(p))\bigr)\\
 &\qquad=\mathbb P\bigl(b(\Phi_i(p))\neq a(\pi(\Phi_i(p)))\bigr)\\
 &\qquad=\mu_\alpha\{b\neq a\circ\pi\}<\eta.
\end{align*}
Tonelli's theorem then gives
\[
 \mathbb E R_i
 =\frac1{V_i}\int_{M_i}
 \mathbb P\bigl(b(\Phi_i(p))\neq a(\Psi_i(p))\bigr)
 \,d\vol_{M_i}(p)<\eta.
\]
Since $R_i\geq0$, Markov's inequality gives
\[
 \mathbb P(R_i>\sqrt\eta)
 \leq\frac{\mathbb E R_i}{\sqrt\eta}<\sqrt\eta.
\]
The primed construction has the same law, so the same argument applies to $R_i'$.
\end{proof}

\begin{lemma}\label{lem:metric-label}
Suppose
$ Z_i^\mathrm{opp}\geq q_*$,
 $R_i,R_i'\leq\sqrt\eta$.
Then
\begin{equation}\label{eq:metric-separation}
 \rho_Y^{M_i}(\Psi_i,\Psi_i')\geq
 \kappa(q_*-2\sqrt\eta)\geq\frac{\kappa q_*}{2}.
\end{equation}
\end{lemma}

\begin{proof}
Let $E_i\subset D_i$ be the set of roots where the two local labels are opposite. Its normalized volume is $Z_i^\mathrm{opp}\geq q_*$. Remove from $E_i$ the roots where either local label differs from the corresponding true factor label. The removed normalized volume is at most $R_i+R_i'\leq2\sqrt\eta$. At every remaining root, one target point lies in $C_0$ and the other in $C_1$. Their $\rho_Y$-distance is at least $\kappa$. Integrating proves the first inequality; the second follows from \eqref{eq:eta-choice}.
\end{proof}

\section{Proof of the main theorem}


We first record the measurability needed for the conditioning argument. Let $\mathscr L_i=L^1(M_i,Y)$ denote the space of measurable maps modulo equality almost everywhere, equipped with the metric $\rho_Y^{M_i}$ (all such maps are integrable because $\rho_Y\leq1$). By \cref{lem:mapavg-borel}, $\mathscr L_i$ is Polish and every averaged model space $\Mapavg(M_i,Y,\rho_Y:L,\delta,\cO)$ is Borel. Joint measurability of $(\omega,p)\mapsto\Psi_i(\omega,p)$ and the standard Fubini-to-$L^1$ measurability theorem imply that $\Psi_i$ is an $\mathscr L_i$-valued random variable. So every conditioning event used below is measurable.

Fix arbitrary microstate parameters: a precompact identity neighborhood $L\subset G$, $\delta>0$, and a weak-star neighborhood $\cO$ of $\nu$. Let
\[
 \mathscr M_i:=\Mapavg(M_i,Y,\rho_Y:L,\delta,\cO).
\]
Define the event
\[
 \mathcal G_i:=\{\Psi_i\in\mathscr M_i\}\cap\{R_i\leq\sqrt\eta\}.
\]
By \cref{cor:good-microstates,lem:Ri},
\begin{equation}\label{eq:good-probability}
 \liminf_i\mathbb P(\mathcal G_i)\geq1-\sqrt\eta.
\end{equation}
Set $c_0:=(1-\sqrt\eta)/2>0$. Then $\mathbb P(\mathcal G_i)\geq c_0$ for all sufficiently large $i$.
For all sufficiently large $i$, let $Q_i$ be the conditional law of $\Psi_i$ given $\mathcal G_i$. It is supported on $\mathscr M_i$.

Set
$
 \eps_0:=\frac{\kappa q_*}{4}>0.
$ We have the following result.

\begin{proposition}\label{prop:collision}
If $\psi,\psi'$ are independent with law $Q_i$, then for all sufficiently large $i$,
\begin{equation}\label{eq:collision}
 \mathbb P\bigl(\rho_Y^{M_i}(\psi,\psi')\leq\eps_0\bigr)
 \leq C_0e^{-\gamma V_i},
\end{equation}
where $C_0=2c_0^{-2}$ and $\gamma>0$ is independent of $L,\delta,\cO$.
\end{proposition}

\begin{proof}
Under the unconditioned product law, if both copies satisfy $R_i,R_i'\leq\sqrt\eta$ and their distance is at most $\eps_0<\kappa q_*/2$, then \cref{lem:metric-label} implies $Z_i^\mathrm{opp}<q_*$. This gives
\[
 \mathbb P\bigl(
 \rho_Y^{M_i}(\Psi_i,\Psi_i')\leq\eps_0,
 \mathcal G_i,\mathcal G_i'
 \bigr)
 \leq\mathbb P(Z_i^\mathrm{opp}<q_*)
 \leq2e^{-\gamma V_i}.
\]
Divide by $\mathbb P(\mathcal G_i)^2\geq c_0^2$. So conditioning on the
good-microstate event changes the collision bound only by a fixed
multiplicative factor, and the exponential decay rate in $V_i$ is preserved.
\end{proof}

\begin{proposition}\label{prop:packing}
For all sufficiently large $i$,
\[
 \Sep_{\eps_0}(\mathscr M_i,\rho_Y^{M_i})
 \geq \left\lfloor\exp\left(\frac\gamma4V_i\right)\right\rfloor.
\]
\end{proposition}

\begin{proof}
Let
$
 N_i:=\left\lfloor e^{\gamma V_i/4}\right\rfloor
$
and sample $\psi_1,\dots,\psi_{N_i}$ independently from $Q_i$. By the union bound and \cref{prop:collision},
\[
 \mathbb P\left(\exists j<k:\rho_Y^{M_i}(\psi_j,\psi_k)\leq\eps_0\right)
 \leq \binom{N_i}{2}C_0e^{-\gamma V_i}
 \leq C_0e^{-\gamma V_i/2}
\]
for all sufficiently large $i$. The right-hand side tends to zero, so there exists an $(\rho_Y^{M_i},\eps_0)$-separated $N_i$-tuple in $\mathscr M_i$.
\end{proof}

Now we can prove  the main theorem.

\begin{proof}[Proof of \cref{thm:main}]
	If $G$ is discrete, then second countability makes $G$ countable and
noncompactness makes it infinite. The coordinate map
$\omega\mapsto(\omega(\{g\}))_{g\in G}$ identifies the Poisson action,
up to the standard inversion conjugacy, with the Bernoulli action having
base law $\operatorname{Pois}(\alpha)$. In this case the local-$G$-space
formalism reduces to the usual finite sofic formalism, so Kerr's theorem
\cite{KerrCPE} gives completely positive entropy for the given
approximation. From now on, assume that $G$ is nondiscrete, which is exactly the
range in which Singh's Theorem~3.4.2 applies.	
	
Let $\pi:(\cN(G),\mu_\alpha)\to(Z,\zeta)$ be a nontrivial measurable factor, and choose a compact metrizable continuous model $(Y,\nu)$ of this factor. The constants $\eps_0>0$ and $\gamma>0$ constructed above depend on the factor, on the chosen finite-window approximation of one nontrivial factor label, and on $\alpha$, but they do not depend on the microstate parameters $L,\delta,\cO$.

For every such triple, \cref{prop:packing} gives
\[
 \limsup_{i\to\infty}\frac1{V_i}
 \log\Sep_{\eps_0}
 \bigl(\Mapavg(M_i,Y,\rho_Y:L,\delta,\cO),\rho_Y^{M_i}\bigr)
 \geq\frac\gamma4.
\]
Taking the infimum over $\cO$, $L$, and $\delta$ gives
\[
 h^{\mathrm{meas,avg}}_{\Sigma,\eps_0}(G,Y,\nu,\rho_Y)
 \geq\frac\gamma4>0.
\]
Because $\eps_0$ and $\gamma$ are independent of the microstate parameters $(L,\delta,\cO)$, the positive lower bound survives the defining infima.

By \cref{prop:avg-standard},
\[
 h^{\mathrm{meas}}_{\Sigma,\eps_0}(G,Y,\nu,\rho_Y)
 =h^{\mathrm{meas,avg}}_{\Sigma,\eps_0}(G,Y,\nu,\rho_Y)>0.
\]
Since the full entropy is the limit as the separation scale decreases,
\[
 h^{\mathrm{meas}}_\Sigma(G,Y,\nu)
 =\lim_{\eps\downarrow0}h^{\mathrm{meas}}_{\Sigma,\eps}(G,Y,\nu,\rho_Y)
 \geq h^{\mathrm{meas}}_{\Sigma,\eps_0}(G,Y,\nu,\rho_Y)>0.
\]
	By \cref{cor:singh-measure-identification}, this is
Singh's standard measure sofic entropy of the factor.  Singh's
measure-conjugacy invariance shows that it is independent of the chosen
compact model. Since the factor was arbitrary, the Poisson action has
completely positive measure sofic entropy relative to $\Sigma$.
\end{proof}

\section{Necessity of the noncompactness hypothesis}

We now show that the noncompactness assumption in \cref{thm:main} is essential for the theorem as stated.

\begin{proposition}\label{prop:compact-obstruction}
Let $G$ be a nontrivial compact lcsc group and let $M_i=G$ with the global left action for every $i$. Then $(M_i)$ is a sofic approximation. For every $\alpha>0$, the Poisson process of intensity $\alpha\Haar$ does not have completely positive measure sofic entropy with respect to this approximation.
\end{proposition}

\begin{proof}
The total-point map
$
 N(\omega)=\omega(G)
$
is a nontrivial invariant factor with law $\lambda=\operatorname{Poi}(\alpha\Haar(G))$ and trivial action. To realize it on a compact target, set
\[
 Y_0:=\{0\}\cup\{y_n:n\in\mathbb N_0\},
 \qquad y_n:=\frac1{n+1},
\]
with the Euclidean metric and the trivial $G$-action. Let $\iota(n)=y_n$ and $\widetilde\lambda=\iota_*\lambda$. Since $\widetilde\lambda(\{0\})=0$, the map $\iota$ is a measure conjugacy from the count-valued factor onto the compact probability system $(Y_0,\widetilde\lambda)$.

Choose distinct $a,b\in\mathbb N_0$ with $\lambda(a),\lambda(b)>0$ and put
$
 \kappa:=|y_a-y_b|>0.
$
The points $y_a,y_b$ are isolated in $Y_0$, so their singleton indicators are continuous. We can choose a weak-star neighborhood $\cO$ of $\widetilde\lambda$ such that every $\beta\in\cO$ satisfies
\[
 \beta(\{y_a\})>\lambda(a)/2,
 \qquad
 \beta(\{y_b\})>\lambda(b)/2.
\]
Let $\overline\Haar:=\Haar/\Haar(G)$ be normalized Haar measure. Let $\varphi:G\to Y_0$ be measurable and write $\beta=\varphi_*\overline\Haar$. Since $(g,p)\mapsto(gp,p)$ preserves $\overline\Haar\times\overline\Haar$, Fubini gives
\begin{align*}
 \frac1{\Haar(G)}\int_G\rho_{Y_0}^{G}(\varphi\circ g,\varphi)\,d\Haar(g)
 &=\iint_{Y_0\times Y_0}|y-z|\,d\beta(y)d\beta(z)\\
 &\geq2\kappa\,\beta(\{y_a\})\beta(\{y_b\})\\
 &>\frac\kappa2\lambda(a)\lambda(b).
\end{align*}
The target action is trivial, so $g\circ\varphi=\varphi$. Taking $L=G$ and any
$
 0<\delta<\frac\kappa2\lambda(a)\lambda(b)
$
shows that
\[
 \Mapavg(G,Y_0,\rho_{Y_0}:G,\delta,\cO)=\varnothing.
\]
Since $\Map\subseteq\Mapavg$, the corresponding standard model space is also empty. So this nontrivial compact factor has entropy $-\infty$, and complete positivity fails for this approximation.
\end{proof}

The proof of \cref{thm:main} also shows exactly where we use noncompactness. By \cref{lem:volume-divergence}, it gives $V_i\to\infty$, which turns a local bounded-difference estimate into an exponential-in-volume packing estimate.

\medskip

\textit{Acknowledgments}.
I am grateful to my advisor, Professor Siming Tu, for introducing me to the subject and his continuous guidance, support and encouragement.



\appendix
\section{Comparison and invariance of spatial entropy }\label{app:singh-comparison}


\subsection{Invariance of spatial entropy}


\begin{lemma}[Measurability of averaged model spaces]\label{lem:mapavg-borel}
	Let $M$ be a finite-volume local $G$-space and let $Y$ be compact metrizable.
	Then $L^1(M,Y)$, equipped with $\rho^M$, is Polish. The empirical-measure
	map
	\[
	\mathsf E_M:L^1(M,Y)\longrightarrow\Prob(Y),
	\qquad \mathsf E_M(\psi)=\psi_*\overline{\vol}_M,
	\]
	is Borel (in fact continuous), the map
	\[
	(g,\psi)\longmapsto\Delta_{\rho,M}(g,\psi)
	\]
	is Borel on $G\times L^1(M,Y)$, and so
	\[
	\psi\longmapsto
	\frac1{\Haar(L)}\int_L\Delta_{\rho,M}(g,\psi)\,d\Haar(g)
	\]
	is Borel for every precompact open subset $L\subset G$. In particular,
	$\Mapavg(M,Y,\rho:L,\delta,\cO)$ is a Borel subset of $L^1(M,Y)$.
\end{lemma}

\begin{proof}
	The Polishness of $L^1(M,Y)$ is standard for a standard probability space
	and a complete separable bounded target. If $\psi_n\to\psi$ in
	$\rho^M$, then $\psi_n\to\psi$ in measure. Uniform continuity of every
	$f\in C(Y)$, followed by bounded convergence in measure, gives
	\[
	\int_M f(\psi_n(p))\,d\overline{\vol}_M(p)
	\longrightarrow
	\int_M f(\psi(p))\,d\overline{\vol}_M(p).
	\]
	So $\mathsf E_M$ is continuous for the weak-star topology on
	$\Prob(Y)$.
	
	Fix $y_*\in Y$. For $g\in G$ define the partial-composition operator
	$C_g:L^1(M,Y)\to L^1(M,Y)$ by
	\[
	(C_g\psi)(p):=
	\begin{cases}
		\psi(g.p),&p\in M[g],\\
		y_*,&p\notin M[g].
	\end{cases}
	\]
	This is well defined on almost-everywhere classes because partial
	translations preserve the canonical measure and so carry null sets to null
	sets. We claim that $(g,\psi)\mapsto C_g\psi$ is Borel. Choose a countable
	dense family $(s_n)$ in $L^1(M,Y)$. For each fixed $n$, the map
	$g\mapsto C_gs_n$ is Borel: the domain of the partial action is Borel, the
	partial action is Borel, and the distance from $C_gs_n$ to any fixed member of
	a countable dense family is a parameterized integral of a bounded Borel
	function. For each $k$, choose Borelly an index $n_k(\psi)$ such that
	\[
	\rho^M(s_{n_k(\psi)},\psi)<2^{-k}.
	\]
	Local measure preservation gives the contraction estimate
	\[
	\rho^M(C_gs_{n_k(\psi)},C_g\psi)
	\leq \rho^M(s_{n_k(\psi)},\psi)<2^{-k}.
	\]
	So $C_gs_{n_k(\psi)}\to C_g\psi$ in $L^1$, proving the claim as a
	pointwise limit of Borel maps.
	
	The map $(g,\psi)\mapsto g\psi$, where $(g\psi)(p)=g(\psi(p))$, is
	continuous from $G\times L^1(M,Y)$ to $L^1(M,Y)$ by joint continuity of the
	action and compactness of $Y$. Also
	$g\mapsto\1_{M[g]}$ is Borel as an $L^1(M)$-valued map, again by
	parameterized integration. This gives
	\begin{align*}
		\Delta_{\rho,M}(g,\psi)
		&=\int_M\1_{M[g]}(p)
		\rho\bigl((C_g\psi)(p),(g\psi)(p)\bigr)
		\,d\overline{\vol}_M(p)\\
		&\qquad+\overline{\vol}_M(M\setminus M[g])
	\end{align*}
	is Borel in $(g,\psi)$, since integration of the pointwise metric is
	continuous in the $L^1$ variables. A final application of parameterized
	integration over $L$ gives the averaged defect, and the claim about
	$\Mapavg$ follows from continuity of $\mathsf E_M$.
\end{proof}


\begin{proposition}[Metric independence]\label{prop:compact-entropy-invariance}
	The quantities
	\[
	h^{\mathrm{meas}}_{\Sigma}(G,Y,\nu,\rho)
	\quad\text{and}\quad
	h^{\mathrm{meas,avg}}_{\Sigma}(G,Y,\nu,\rho)
	\]
	are independent of the compatible metric $\rho$ on the fixed compact model $Y$.
\end{proposition}

\begin{proof}
	It is enough, by \cref{prop:avg-standard}, to prove the claim for the averaged entropy. Let $\rho$ and $\rho'$ be compatible metrics bounded by $1$. Fix a separation scale $r\in(0,1)$. By uniform equivalence on the compact space $Y$, choose $u\in(0,1)$ such that
	\begin{equation}\label{eq:separation-modulus}
		\rho'(y,z)<u\Longrightarrow \rho(y,z)<r/2.
	\end{equation}
	Put $t:=ur/4$. If $(\rho')^M(\psi,\psi')\leq t$, then Markov's
inequality shows that the set where $\rho'\geq u$ has normalized volume
at most $r/4$. On its complement, \eqref{eq:separation-modulus} gives
$\rho<r/2$. Since $\rho\leq1$,
\[
\rho^M(\psi,\psi')<\frac r2+\frac r4=\frac{3r}{4}<r.
\]
By contraposition, $\rho^M(\psi,\psi')>r$ implies
$(\rho')^M(\psi,\psi')>t$. So every
$(\rho^M,r)$-separated family is $(\rho'^M,t)$-separated, with the
strict inequality required by the definition of separation.
	
	Now fix an averaged-equivariance tolerance $\delta'>0$. Choose $a\in(0,\delta'/2)$ and then choose $s\in(0,1)$ such that
	\begin{equation}\label{eq:equivariance-modulus}
		\rho(y,z)<s\Longrightarrow \rho'(y,z)<a.
	\end{equation}
	Splitting the defined part of the domain into $\{\rho<s\}$ and its complement, and using $s\leq1$ for the undefined-domain penalty, gives
	\begin{equation}\label{eq:defect-comparison}
		\Delta_{\rho',M}(g,\psi)
		\leq a+s^{-1}\Delta_{\rho,M}(g,\psi).
	\end{equation}
	Set $\delta:=s(\delta'-a)>0$. Then
	\begin{equation}\label{eq:map-inclusion}
		\Mapavg(M,Y,\rho:L,\delta,\cO)
		\subseteq
		\Mapavg(M,Y,\rho':L,\delta',\cO).
	\end{equation}
	Combining separation with \eqref{eq:map-inclusion}, for every $L$, $\cO$, and $\delta'>0$ we get
	\begin{align*}
		&\limsup_i\frac1{V_i}
		\log\Sep_t\bigl(\Mapavg(M_i,Y,\rho':L,\delta',\cO),\rho'^{M_i}\bigr)\\
		&\qquad\geq
		\limsup_i\frac1{V_i}
		\log\Sep_r\bigl(\Mapavg(M_i,Y,\rho:L,\delta,\cO),\rho^{M_i}\bigr).
	\end{align*}
	Taking the infima over $\delta'$, $L$, and $\cO$, and then letting $r\downarrow0$, gives
	\[
	h^{\mathrm{meas,avg}}_{\Sigma}(G,Y,\nu,\rho')
	\geq h^{\mathrm{meas,avg}}_{\Sigma}(G,Y,\nu,\rho).
	\]
	Interchanging $\rho$ and $\rho'$ proves equality, and \cref{prop:avg-standard} transfers it to the pointwise measure sofic entropy.
\end{proof}

\begin{proposition}[Topological-conjugacy invariance]\label{prop:top-conj-invariance}
	Let $G\curvearrowright(Y,\nu)$ and $G\curvearrowright(Y',\nu')$ be compact metrizable continuous pmp systems. Suppose that
	\[
	\Phi:Y\longrightarrow Y'
	\]
	is a $G$-equivariant homeomorphism with $\Phi_*\nu=\nu'$. Then
	\[
	h^{\mathrm{meas}}_\Sigma(G,Y,\nu)=h^{\mathrm{meas}}_\Sigma(G,Y',\nu')
	\]
	and
	\[
	h^{\mathrm{meas,avg}}_\Sigma(G,Y,\nu)
	=h^{\mathrm{meas,avg}}_\Sigma(G,Y',\nu').
	\]
	So the spatial formula is invariant under measure-preserving topological $G$-conjugacy.
\end{proposition}

\begin{proof}
	Let $\rho'$ be a compatible metric on $Y'$ bounded by $1$, and pull it back to the compatible metric
	\[
	\widetilde\rho(y,z):=\rho'(\Phi(y),\Phi(z))
	\qquad(y,z\in Y)
	\]
	on $Y$. For every finite-volume local $G$-space $M$, composition with $\Phi$ gives a bijection
	\[
	T_M:Y^M\longrightarrow (Y')^M,
	\qquad T_M(\psi):=\Phi\circ\psi.
	\]
	Because $\Phi$ is $G$-equivariant, for every $g\in G$ and every measurable $\psi:M\to Y$ we have
	\begin{align*}
		\Delta_{\rho',M}(g,T_M\psi)
		&=\int_{M[g]}\rho'\bigl(\Phi(\psi(g.p)),g\Phi(\psi(p))\bigr)
		\,d\overline{\vol}_M(p)
		+\overline{\vol}_M(M\setminus M[g])\\
		&=\int_{M[g]}\widetilde\rho\bigl(\psi(g.p),g\psi(p)\bigr)
		\,d\overline{\vol}_M(p)
		+\overline{\vol}_M(M\setminus M[g])\\
		&=\Delta_{\widetilde\rho,M}(g,\psi).
	\end{align*}
	Also,
	\[
	(T_M\psi)_*\overline{\vol}_M
	=\Phi_*\bigl(\psi_*\overline{\vol}_M\bigr),
	\]
	and for all $\psi,\psi':M\to Y$,
	\[
	(\rho')^M(T_M\psi,T_M\psi')
	=\widetilde\rho^M(\psi,\psi').
	\]
	
	The pushforward map
	\[
	\Phi_*:\operatorname{Prob}(Y)\longrightarrow\operatorname{Prob}(Y')
	\]
	is a homeomorphism for the weak-star topologies, with inverse $(\Phi^{-1})_*$. So if $\cO'$ is a weak-star neighborhood of $\nu'$ and
	\[
	\cO:=\Phi_*^{-1}(\cO'),
	\]
	then $\cO$ is a weak-star neighborhood of $\nu$, and for every precompact open identity neighborhood $L\subset G$ and every $\delta>0$,
	\begin{align*}
		T_M\bigl(\Map(M,Y,\widetilde\rho:L,\delta,\cO)\bigr)
		&=\Map(M,Y',\rho':L,\delta,\cO'),\\
		T_M\bigl(\Mapavg(M,Y,\widetilde\rho:L,\delta,\cO)\bigr)
		&=\Mapavg(M,Y',\rho':L,\delta,\cO').
	\end{align*}
	These bijections are isometries for the normalized integral metrics. So for every $\eps>0$, every $i$, and the corresponding neighborhoods $\cO,\cO'$,
	\begin{align*}
		\Sep_\eps\bigl(\Map(M_i,Y,\widetilde\rho:L,\delta,\cO),
		\widetilde\rho^{M_i}\bigr)=
		\Sep_\eps\bigl(\Map(M_i,Y',\rho':L,\delta,\cO'),
		(\rho')^{M_i}\bigr),
	\end{align*}
	and the same equality holds for the averaged model spaces. Since $\Phi_*$ bijects the weak-star neighborhood filters of $\nu$ and $\nu'$, taking the infima in the definitions gives
	\[
	h^{\mathrm{meas}}_{\Sigma,\eps}(G,Y,\nu,\widetilde\rho)
	=h^{\mathrm{meas}}_{\Sigma,\eps}(G,Y',\nu',\rho')
	\]
	for every $\eps>0$, and also for $h^{\mathrm{meas,avg}}_{\Sigma,\eps}$. Letting $\eps\downarrow0$ gives equality of the corresponding entropies for the metrics $\widetilde\rho$ and $\rho'$. Finally, metric independence from \cref{prop:compact-entropy-invariance} allows $\widetilde\rho$ and $\rho'$ to be replaced by arbitrary compatible metrics on $Y$ and $Y'$, respectively.
\end{proof}
	\begin{remark}\label{rem:top-vs-measure-conjugacy}

		The previous proposition compares compact models only when they are
		topologically conjugate. Two compact models of the same standard pmp action
		need only be measure conjugate. The required compact-model independence is
		proved below by identifying the local formula with Singh's
		measure-conjugacy invariant entropy.

\end{remark}


\subsection{Comparison with Singh’s pseudometric entropy}
In this subsection we verify, in the common compact-model setting, that the
spatial measure entropy of \cref{def:spatial-entropy} agrees with the
pseudometric microstate entropy introduced by Singh in
\cite[Definitions~9--11]{SinghThesis}. The two formulations differ in three ways. First, the present paper uses local $G$-spaces, so the translations are only
partially defined. Singh uses globally defined measurable maps. Second, the
present paper uses an $L^1$ model pseudometric, while Singh
uses the corresponding $L^2$ pseudometric. Third, the empirical-distribution
condition is expressed here by a weak-star neighborhood and in Singh's definition
by finitely many continuous test functions. We show that none of these differences affects the entropy.

Let $G\curvearrowright (X,\mu)$ be a jointly continuous
pmp action on a compact metrizable space  and $\Sigma=(M_i)_{i\geq 1}$ be a sofic approximation
in the sense of Definition \ref{def:sofic-approximation}, and let
$\rho:X\times X\to[0,1]$ be a continuous generating pseudometric.

	Locally compact sofic groups are unimodular. So the right Haar measure
used to define the canonical measures on the local $G$-spaces is also a left
Haar measure and may be used in Singh's Definition~4 without changing the
normalization. The completion constructed below depends only on the original
approximation $\Sigma$, not on the action, factor, or compact model.

\subsubsection{From local \texorpdfstring{$G$}{G}-spaces to Singh approximations}

Recall that for $g\in G$,
\[
M[g]:=\{p\in M:(g,p)\text{ lies in the domain of the partial action}\}.
\]
Choose a Borel completion of the partial action by setting
\begin{equation}\label{eq:completion}
	\widetilde L_i(g,p):=
	\begin{cases}
		g.p, & p\in M[g],\\
		p,   & p\notin M[g].
	\end{cases}
\end{equation}
The map $\widetilde L_i:G\times M_i\to M_i$ is Borel because the domain of the
partial action is open and the partial action is continuous there.

\begin{lemma}\label{lem:completion-is-singh}
	The sequence
$
	\widetilde\Sigma
	:=(M_i,\vol_i,\widetilde L_i)_{i\geq 1}
$
	is a sofic approximation in the sense of Singh.
\end{lemma}

\begin{proof}
	Fix a precompact open neighborhood $U$ of the identity. Let $M_i[U]$ be as in
	\cref{def:local-sofic-model} and put
	\[
	\beta_i(U):=\overline\vol_i(M_i\setminus M_i[U]).
	\]
	The definition of a sofic approximation implies
	\begin{equation}\label{eq:beta-to-zero}
		\beta_i(U)\longrightarrow 0.
	\end{equation}
	In fact, for all sufficiently large $i$, one has $U\subset U_i$ and so
	$M_i[U_i]\subset M_i[U]$.
	
	If $p\in M_i[U]$, then $g\mapsto g.p$ is a homeomorphism from $U$ onto an open
	neighborhood of $p$, and
	\[
	g.h.p=gh.p
	\qquad(g,h,gh\in U).
	\]
	Also, the defining property of the canonical measure gives
	\[
	\vol_i(\{g.p:g\in A\})=\Haar(A)
	\]
	for every Borel set $A\subset U$. So the restriction of
	$g\mapsto\widetilde L_i(g,p)=g.p$ to $U$ is a measure-space isomorphism onto its
	image, it sends the identity to $p$, and it satisfies the required local
	multiplication law. These are exactly conditions~(1)--(3) of
	\cite[Definition~4]{SinghThesis}. Since $M_i[U]$ has normalized measure tending to one by
	\eqref{eq:beta-to-zero}, the sequence $\widetilde\Sigma$ satisfies Singh's
	sofic-approximation axioms.
\end{proof}








\subsubsection{The two microstate spaces}



For measurable maps $\phi,\psi:M_i\to X$, define
\begin{align*}
	\rho_{1,i}(\phi,\psi)
	&:=\int_{M_i}\rho(\phi(p),\psi(p))\,d\overline\vol_i(p),\\
	\rho_{2,i}(\phi,\psi)
	&:=\left(\int_{M_i}\rho(\phi(p),\psi(p))^2
	\,d\overline\vol_i(p)\right)^{1/2}.
\end{align*}
For $g\in G$, define the Bowen equivariance error
\begin{equation}\label{eq:bowen-error}
E^B_i(\phi,g)
	:=\overline\vol_i(M_i\setminus M_i[g])
	+\int_{M_i[g]}
	\rho\bigl(\phi(g.p),g\phi(p)\bigr)\,d\overline\vol_i(p),
\end{equation}
and the completed Singh equivariance error
\begin{equation}\label{eq:singh-error}
E^S_i(\phi,g)
	:=\left(\int_{M_i}
	\rho\bigl(\phi(\widetilde L_i(g,p)),g\phi(p)\bigr)^2
	\,d\overline\vol_i(p)\right)^{1/2}.
\end{equation}
So $E^B_i(\phi,g)$ is exactly the defect appearing in \eqref{eq:defect}.

If $\mathcal F\subset C(X)$ is finite and $\eta>0$, let
\begin{equation}\label{eq:weak-star-basic}
	\cO(\mathcal F,\eta)
	:=\left\{\nu\in\Prob(X):
	\left|\int f\,d\nu-\int f\,d\mu\right|<\eta
	\text{ for every }f\in\mathcal F\right\}.
\end{equation}
These sets form a neighborhood basis of $\mu$ for the weak-star topology on
$\Prob(X)$.

For a precompact open neighborhood $U\subset G$ and $\delta>0$, write
\begin{align*}
	\Map_B(i;U,\delta,\mathcal F,\eta)
	:=\{\phi:M_i\to X:\;&E^B_i(\phi,g)<\delta\text{ for all }g\in U,\\
	&\phi_*\overline\vol_i\in\cO(\mathcal F,\eta)\},
\end{align*}
and
\begin{align*}
	\Map_S(i;U,\delta,\mathcal F,\eta)
	:=\{\phi:M_i\to X:\;&E^S_i(\phi,g)<\delta\text{ for all }g\in U,\\
	&\phi_*\overline\vol_i\in\cO(\mathcal F,\eta)\}.
\end{align*}
The first is the pointwise microstate space from Section~2, restricted to the basic
weak-star neighborhood \eqref{eq:weak-star-basic}. The second is Singh's spatial
microstate space, except that we temporarily allow separate tolerances for
equivariance and empirical distribution. This does not alter Singh's entropy:
the diagonal family formed by setting both tolerances equal is cofinal, because
for any $\delta,\eta>0$ one may replace them by $\min\{\delta,\eta\}$.

\begin{lemma}\label{lem:error-comparison}
	For every precompact open $U\subset G$, every $g\in U$, every $i$, and every
	measurable $\phi:M_i\to X$,
	\begin{equation}\label{eq:error-comparison}
		E^S_i(\phi,g)^2\le E^B_i(\phi,g)
		\quad\text{and}\quad
		E^B_i(\phi,g)\le E^S_i(\phi,g)+\beta_i(U).
	\end{equation}
	These inequalities give, for every $0<\delta<1$,
	\begin{equation}\label{eq:map-inclusions}
		\Map_B(i;U,\delta^2,\mathcal F,\eta)
		\subseteq
		\Map_S(i;U,\delta,\mathcal F,\eta),
	\end{equation}
	and, for all sufficiently large $i$,
	\begin{equation}\label{eq:map-inclusions-reverse}
		\Map_S(i;U,\delta/2,\mathcal F,\eta)
		\subseteq
		\Map_B(i;U,\delta,\mathcal F,\eta).
	\end{equation}
\end{lemma}

\begin{proof}
	Fix $g\in U$. Since $M_i[U]\subset M_i[g]$, one has
	\begin{equation}\label{eq:missing-mass}
		\overline\vol_i(M_i\setminus M_i[g])\le\beta_i(U).
	\end{equation}
	Put
	\[
	a(p):=\rho(\phi(g.p),g\phi(p))\quad(p\in M_i[g])
	\]
	and
	\[
	c(p):=\rho(\phi(\widetilde L_i(g,p)),g\phi(p))
	\quad(p\in M_i\setminus M_i[g]).
	\]
	Because $0\le a,c\le1$ and $\widetilde L_i(g,p)=g.p$ on $M_i[g]$,
	\begin{align*}
		E^S_i(\phi,g)^2
		&=\int_{M_i[g]} a(p)^2\,d\overline\vol_i(p)
		+\int_{M_i\setminus M_i[g]}c(p)^2\,d\overline\vol_i(p)\\
		&\le\int_{M_i[g]} a(p)\,d\overline\vol_i(p)
		+\overline\vol_i(M_i\setminus M_i[g])
		=E^B_i(\phi,g).
	\end{align*}
	But by Cauchy--Schwarz and \eqref{eq:missing-mass},
	\begin{align*}
		E^B_i(\phi,g)
		&\le\beta_i(U)+\int_{M_i[g]} a(p)\,d\overline\vol_i(p)\\
		&\le\beta_i(U)
		+\left(\int_{M_i[g]} a(p)^2\,d\overline\vol_i(p)\right)^{1/2}
		\le\beta_i(U)+E^S_i(\phi,g).
	\end{align*}
	This proves \eqref{eq:error-comparison}. The inclusion
	\eqref{eq:map-inclusions} follows directly. By
	\eqref{eq:beta-to-zero}, for all sufficiently large $i$ one has
	$\beta_i(U)<\delta/2$, and then \eqref{eq:map-inclusions-reverse} follows from
	the second inequality in \eqref{eq:error-comparison}.
\end{proof}

\begin{lemma}\label{lem:l1-l2-separation}
	For all measurable $\phi,\psi:M_i\to X$,
	\begin{equation}\label{eq:l1-l2-model}
		\rho_{1,i}(\phi,\psi)
		\le\rho_{2,i}(\phi,\psi)
		\le\rho_{1,i}(\phi,\psi)^{1/2}.
	\end{equation}
	These inequalities give, for every $A\subset\Map(M_i,X)$ and every $0<\eps<1$,
	\begin{equation}\label{eq:sep-comparison}
		\Sep_\eps(A,\rho_{1,i})
		\le\Sep_\eps(A,\rho_{2,i})
		\le\Sep_{\eps^2}(A,\rho_{1,i}).
	\end{equation}
\end{lemma}

\begin{proof}
	Apply $\|f\|_1\le\|f\|_2$ and $f^2\le f$ to the function
	$f(p)=\rho(\phi(p),\psi(p))\in[0,1]$. The separation-number inequalities are
	immediate from \eqref{eq:l1-l2-model}.
\end{proof}

\subsubsection{Equality of the entropies}

For $0<\eps<1$, define
\begin{align*}
	b_\eps
	:=\inf_{\mathcal F}\inf_{\eta>0}\inf_U\inf_{\delta>0}
	\limsup_{i\to\infty}\frac{1}{\vol_i(M_i)}
	\log\Sep_\eps
	\bigl(\Map_B(i;U,\delta,\mathcal F,\eta),\rho_{1,i}\bigr),
\end{align*}
where $\mathcal F$ ranges over finite subsets of $C(X)$ and $U$ over precompact
open neighborhoods of the identity. Because the sets
$\cO(\mathcal F,\eta)$ form a weak-star neighborhood basis,
\begin{equation}\label{eq:bowen-b-epsilon}
	h^{\mathrm{meas}}_\Sigma(G,X,\mu,\rho)=\sup_{0<\eps<1} b_\eps.
\end{equation}
In the same way, let
\begin{align*}
	s_\eps
	:=\inf_{\mathcal F}\inf_{\eta>0}\inf_U\inf_{\delta>0}
	\limsup_{i\to\infty}\frac{1}{\vol_i(M_i)}
	\log\Sep_\eps
	\bigl(\Map_S(i;U,\delta,\mathcal F,\eta),\rho_{2,i}\bigr).
\end{align*}
By the cofinality observation above,
\[
\sup_{0<\eps<1}s_\eps
=h_{\widetilde\Sigma}^{\mathrm{Singh,pseudo}}(\mu,\rho),
\]
where the right-hand side is Singh's pseudometric entropy computed with the
completed approximation $\widetilde\Sigma$.

	\begin{theorem}[Singh--Bowen comparison]\label{thm:identification-bowen}
	Under the assumptions of this subsection,
	\begin{equation}\label{eq:main-equivalence}
		h^{\mathrm{meas}}_\Sigma(G,X,\mu,\rho)
		=h_{\widetilde\Sigma}^{\mathrm{Singh,pseudo}}(\mu,\rho).
	\end{equation}
	Here $\widetilde\Sigma$ is the Singh approximation built by completing
	the same local approximation $\Sigma$. The common value is independent of
	how the partial action is completed outside its domain, and the completed
	approximation depends only on $\Sigma$, not on the action or compact model.
\end{theorem}

\begin{proof}
	Fix $0<\eps<1$. By \eqref{eq:map-inclusions} and the first inequality in
	\eqref{eq:sep-comparison}, for every choice of
	$U,\mathcal F,\eta,$ and $0<\delta<1$,
	\begin{align*}
		\Sep_\eps
		\bigl(\Map_B(i;U,\delta^2,\mathcal F,\eta),\rho_{1,i}\bigr)\le
		\Sep_\eps
		\bigl(\Map_S(i;U,\delta,\mathcal F,\eta),\rho_{2,i}\bigr).
	\end{align*}
	Taking normalized logarithms, the limsup, and then the indicated infima gives
	\begin{equation}\label{eq:b-less-s}
		b_\eps\le s_\eps.
	\end{equation}
	Here replacing $\delta$ by $\delta^2$ does not affect the infimum, since
	$\{\delta^2:0<\delta<1\}$ is cofinal at zero.

	For the reverse direction, by \eqref{eq:map-inclusions-reverse}, for all sufficiently large $i$,
	\[
	\Map_S(i;U,\delta/2,\mathcal F,\eta)
	\subseteq\Map_B(i;U,\delta,\mathcal F,\eta).
	\]
	Using the second inequality in \eqref{eq:sep-comparison},
	\begin{align*}
		\Sep_\eps
		\bigl(\Map_S(i;U,\delta/2,\mathcal F,\eta),\rho_{2,i}\bigr)
		\le
		\Sep_{\eps^2}
		\bigl(\Map_B(i;U,\delta,\mathcal F,\eta),\rho_{1,i}\bigr)
	\end{align*}
	for all sufficiently large $i$. This gives
	\begin{equation}\label{eq:s-less-b}
		s_\eps\le b_{\eps^2}.
	\end{equation}
	Combining \eqref{eq:b-less-s} and \eqref{eq:s-less-b},
	\[
	b_\eps\le s_\eps\le b_{\eps^2}.
	\]
	Taking the supremum over $0<\eps<1$ gives
	\[
	\sup_{0<\eps<1}b_\eps
	=\sup_{0<\eps<1}s_\eps,
	\]
	because $\eps\mapsto\eps^2$ maps $(0,1)$ onto $(0,1)$. This proves
	\eqref{eq:main-equivalence}. The argument used no property of the completion
	outside the estimate $0\le c\le1$; so every Borel completion gives the same
	entropy.
\end{proof}

	\begin{corollary}[Identification with Singh's measure entropy]
	\label{cor:singh-measure-identification}
	Assume in addition that $G$ is nondiscrete. In the notation of Singh's
	thesis, for every dynamically generating sequence $P$, every countable set
	$\Gamma\subset G$ which generates $P$ in Singh's sense, and every
	constant $C\geq1$,
	\begin{align*}
		h^{\mathrm{meas}}_\Sigma(G,X,\mu,\rho)
		&=h_{\widetilde\Sigma}^{\mathrm{Singh,pseudo}}(\mu,\rho)\\
		&=h_{\widetilde\Sigma}^{\mathrm{meas,Singh}}
		\bigl(G\curvearrowright(X,\mu)\bigr)\\
		&=h_{\operatorname{Hom}_C}(P,\Gamma)
		=h_{\operatorname{UP}_C}(P,\Gamma).
	\end{align*}
	So the local spatial entropy is invariant under measure
	conjugacy and is independent of the compact metrizable continuous model and
	compatible metric used to represent the underlying standard pmp action.
\end{corollary}

\begin{proof}
	The first equality is \cref{thm:identification-bowen}. Singh's
	Theorem~3.4.2 identifies the unrestricted and $L^2$-bounded pseudometric
	formulas for nondiscrete groups. Proposition~3.4.3 identifies the latter
	with the unital-homomorphism formula, while Theorem~3.3.2 and
	Definition~11 identify the homomorphism and unital-positive-map formulas and
	define their common value to be the measure entropy of the action. The
	independence of the generating data proved in Chapter~3 makes this common
	value a measure-conjugacy invariant; see
	\cite[Theorems~3.2.3, 3.3.2, 3.4.2, Proposition~3.4.3, and
	Definition~11]{SinghThesis}.
\end{proof}


\begin{remark}\label{rem:scope-comparison}

	\Cref{thm:identification-bowen} compares the local formula for
	$\Sigma$ with Singh's formula for the particular completion
	$\widetilde\Sigma$ built from $\Sigma$; it does not compare entropies
	attached to two unrelated sofic approximation sequences.  Compactness of the
	target is used to work directly in Singh's spatial framework.  Since every
	standard pmp action of an lcsc group has a compact metrizable continuous
	model, no noncompact-target comparison is needed for the present theorem.

\end{remark}





\begin{thebibliography}{99}
	\bibitem{AustinBurton}
	T.~Austin and P.~Burton,
	\emph{Uniform mixing and completely positive sofic entropy},
	J. Anal. Math. \textbf{138} (2019), 597--612.
	\bibitem{BachmannPeccati}
	S.~Bachmann and G.~Peccati,
	\emph{Concentration bounds for geometric Poisson functionals: logarithmic Sobolev inequalities revisited},
	Electron. J. Probab. \textbf{21} (2016), Paper No. 6, 44 pp.
	
	\bibitem{BowenSofic}
	L.~Bowen,
	\emph{Measure conjugacy invariants for actions of countable sofic groups},
	J. Amer. Math. Soc. \textbf{23} (2010), 217--245.
	\bibitem{B12}L. Bowen,
	\textit{Sofic entropy and amenable groups}, Ergod. Th. Dynam. Sys. \textbf{32} (2012), 427--466.
	
	\bibitem{B20}L. Bowen,
	\textit{Examples in the entropy theory of countable group actions}, Ergod. Th. Dynam. Sys. \textbf{40} (2020), 2593--2680.
	
	\bibitem{BowenLCEntropy}
	L.~Bowen,
	\emph{Locally compact sofic entropy theory},
	Trans. Amer. Math. Soc. Ser. B \textbf{12} (2025), 165--190.
	\bibitem{BowenBurton}
	L.~Bowen and P.~Burton,
	\emph{Locally compact sofic groups},
	Israel J. Math. \textbf{251} (2022), 239--270.
	\bibitem{Glasner}
	E.~Glasner,
	\emph{Ergodic Theory via Joinings},
	Mathematical Surveys and Monographs, vol.~101,
	American Mathematical Society, Providence, RI, 2003.

	
	
	
	

	

	
	\bibitem{KerrCPE}
	D.~Kerr,
	\emph{Bernoulli actions of sofic groups have completely positive entropy},
	Israel J. Math. \textbf{202} (2014), 461--474.
	\bibitem{KerrLi1} D. Kerr and H. Li, \textit{Entropy and the variational principle for actions of sofic groups}, Invent. Math. \textbf{186} (2011), 501--558.
	
	\bibitem{KerrLi2} D. Kerr and H. Li, \textit{Soficity, amenability, and dynamical entropy}, Amer. J. Math. \textbf{135} (2013), 721-761.
		
	\bibitem{OrnsteinWeiss}
	D.~S.~Ornstein and B.~Weiss,
	\emph{Entropy and isomorphism theorems for actions of amenable groups},
	J. Analyse Math. \textbf{48} (1987), 1--141.
	
	\bibitem{RudolphWeiss}
	D.~J.~Rudolph and B.~Weiss,
	\emph{Entropy and mixing for amenable group actions},
	Ann. of Math. (2) \textbf{151} (2000), 1119--1150.
	\bibitem{SinghThesis}
	S.~Singh,
	\emph{Entropy theory for locally compact sofic groups},
	Ph.D. thesis, The University of Texas at Austin, 2016.
	
	\bibitem{Var63}
	V. S. Varadarajan, \emph{Groups of automorphisms of Borel spaces}, Trans.
	Amer. Math. Soc. \textbf{109} (1963), 191--220.
	
	\bibitem{Wu}
	L.~Wu,
	\emph{A new modified logarithmic Sobolev inequality for Poisson point processes and several applications},
	Probab. Theory Related Fields \textbf{118} (2000), 427--438.
	
\end{thebibliography}

\end{document}
